\documentclass[11pt,letterpaper]{article}

\usepackage[spanish,es-nosectiondot]{babel}
\usepackage[utf8]{inputenc}
\usepackage[T1]{fontenc}
\usepackage{amsmath}
\usepackage{booktabs}
\usepackage{geometry}
\usepackage{longtable}
\usepackage{pdflscape}
\usepackage{threeparttable}
\usepackage{xurl}

\geometry{margin=2.5cm}

\title{Reporte del modelo logit multinomial territorial}
\author{Juan Vicente Onetto Romero}
\date{Mayo 2026}

\begin{document}

\maketitle

\tableofcontents

\section*{Nota inicial}

Este documento corresponde a un reporte de trabajo sobre la especificación, estimación e
interpretación del modelo logit multinomial territorial desarrollado hasta esta etapa de la
investigación. Su objetivo es ordenar la formulación del modelo, las variables utilizadas, los
resultados principales y las decisiones metodológicas pendientes, más que presentar una versión
definitiva del modelo final de tesis.

\section{Descripción del problema y alternativas}

Este reporte documenta un modelo logit multinomial orientado a estudiar la elección del medio de pago en viajes de transporte público en Santiago. En particular, el modelo distingue entre viajes pagados con tarjeta BIP y viajes pagados mediante códigos QR, separando estos últimos según el canal de pago utilizado. El objetivo es analizar cómo atributos del viaje, condiciones operacionales del sistema y características territoriales del entorno de origen se asocian con la utilidad relativa de cada alternativa de pago.

La unidad de observación corresponde a un viaje individual observado en la muestra de estimación. Cada viaje se vincula a una zona de inicio \texttt{ZONA777}, definida a partir del paradero o estación donde comienza el viaje. Las \texttt{ZONA777} corresponden a la zonificación espacial utilizada para agregar viajes, oferta y atributos territoriales del sistema de transporte de Santiago. En este reporte se usan como unidad espacial de referencia para incorporar al modelo información contextual de la zona donde se inicia el viaje, incluyendo variables de demanda, oferta, características sociodemográficas y entorno construido.

El conjunto de alternativas del modelo está compuesto por tres categorías mutuamente excluyentes:
\begin{itemize}
    \item \texttt{BIP}: viaje pagado con tarjeta BIP.
    \item \texttt{QR\_RED}: viaje pagado mediante QR asociado a la aplicación o canal oficial de Red Metropolitana de Movilidad.
    \item \texttt{QR\_OTHER}: viaje pagado mediante otros canales o aplicaciones QR distintos de \texttt{QR\_RED}.
\end{itemize}

La separación entre \texttt{QR\_RED} y \texttt{QR\_OTHER} responde a que el pago QR no constituye necesariamente una alternativa homogénea. Distintos canales QR pueden estar asociados a mecanismos de adopción, disponibilidad tecnológica, patrones de uso y contextos territoriales diferentes. Por esta razón, el modelo no se limita a comparar tarjeta BIP versus QR agregado, sino que permite estimar diferencias en la utilidad relativa de cada tipo de pago digital frente al uso de BIP.

\section{Especificación MNL y función de utilidad}

\subsection{Marco general de utilidad aleatoria}

La especificación del modelo sigue el marco de utilidad aleatoria, ampliamente utilizado en modelos de elección discreta en transporte. Bajo este enfoque, se asume que cada viaje \(n\) tiene asociada una utilidad latente para cada alternativa de medio de pago \(i\). La alternativa observada corresponde a aquella que entrega la mayor utilidad dentro del conjunto de alternativas disponibles.

En este caso, el conjunto de elección es fijo y está compuesto por tres medios de pago:

\begin{equation}
    C_n = \{\texttt{BIP}, \texttt{QR\_RED}, \texttt{QR\_OTHER}\}.
\end{equation}

Para cada viaje \(n\) y alternativa \(i \in C_n\), la utilidad total se define como:

\begin{equation}
    U_{ni} = V_{ni} + \varepsilon_{ni},
\end{equation}

donde \(U_{ni}\) corresponde a la utilidad total de la alternativa \(i\) para el viaje \(n\), \(V_{ni}\) es la utilidad sistemática u observable, y \(\varepsilon_{ni}\) es un componente no observado por el analista. Este último término captura factores no incluidos explícitamente en el modelo, como preferencias individuales no observadas, condiciones específicas del viaje o errores de medición.

La utilidad sistemática \(V_{ni}\) se especifica como una combinación lineal de dos tipos de
variables: atributos que pueden variar entre alternativas dentro de un mismo viaje y variables
comunes al viaje o a la zona de origen. En forma general:

\begin{equation}
    V_{ni}
    =
    ASC_i
    +
    \sum_{k \in K^{AS}} \beta^{AS}_{ki} X_{kni}
    +
    \sum_{m \in K^{C}} \beta^{C}_{mi} Z_{mn}.
\end{equation}

En esta expresión, \(ASC_i\) corresponde a la constante específica de la alternativa \(i\).
El término \(X_{kni}\) representa variables alternativa-específicas, cuyo valor puede cambiar entre
alternativas para un mismo viaje \(n\), como tiempos de viaje, esperas y transbordos. Por su parte,
\(Z_{mn}\) representa variables comunes al viaje o a la zona de origen, como año, franja horaria,
demanda, oferta, variables censales y variables OSM. Aunque estas últimas tienen el mismo valor para
las tres alternativas de una observación, sus coeficientes pueden diferir por alternativa,
permitiendo capturar cómo el contexto temporal o territorial modifica la utilidad relativa de cada
medio de pago.

La constante específica de alternativa captura la utilidad media de una alternativa que no es explicada por las variables observadas. En el modelo, \texttt{BIP} se utiliza como alternativa base de normalización para las constantes:

\begin{equation}
    ASC_{\texttt{BIP}} = 0.
\end{equation}

En consecuencia, las constantes estimadas para \texttt{QR\_RED} y \texttt{QR\_OTHER} se interpretan respecto de \texttt{BIP}. Por ejemplo, si \(ASC_{\texttt{QR\_RED}}\) es negativo, esto indica que, manteniendo constantes las demás variables, \texttt{QR\_RED} presenta una menor utilidad media relativa que \texttt{BIP}.

Bajo el supuesto de que los términos no observados \(\varepsilon_{ni}\) son independientes e idénticamente distribuidos Gumbel tipo I, la probabilidad de que el viaje \(n\) utilice la alternativa \(i\) toma la forma logit multinomial:

\begin{equation}
    P_{ni}
    =
    \frac{\exp(V_{ni})}
    {\sum_{j \in C_n} \exp(V_{nj})}.
\end{equation}

Así, el modelo transforma las utilidades sistemáticas de \texttt{BIP}, \texttt{QR\_RED} y \texttt{QR\_OTHER} en probabilidades de elección. Por esta razón, los coeficientes deben interpretarse siempre en términos relativos entre alternativas: un coeficiente positivo aumenta la utilidad sistemática de una alternativa respecto de las demás, mientras que un coeficiente negativo la reduce, manteniendo constantes el resto de las variables del modelo.

\subsection{Duda metodológica pendiente: Variables alternativas-específicas y variables comunes}

La especificación combina dos tipos de variables. El primer grupo corresponde a variables que pueden variar entre alternativas dentro de una misma observación, como el tiempo en vehículo, los tiempos de espera y el número de transbordos asociados a cada alternativa. Para estas variables, el coeficiente de \texttt{BIP} es directamente identificable, ya que el valor de la variable puede diferir entre \texttt{BIP}, \texttt{QR\_RED} y \texttt{QR\_OTHER}.

El segundo grupo corresponde a variables comunes a las alternativas dentro de una misma observación. En este grupo se incluyen variables temporales y de contexto, como el año, la franja horaria, la demanda y oferta del entorno, además de variables territoriales provenientes del Censo y de OpenStreetMap. Estas variables tienen el mismo valor para las tres alternativas de una observación. Por esta razón, no se identifican tres efectos absolutos independientes, sino diferencias relativas entre alternativas.

Esto puede verse directamente en la probabilidad logit. Si una misma variable \(X_{kn}\) entra con el mismo valor en todas las alternativas, sumar un término común \(c X_{kn}\) a todas las utilidades no modifica las probabilidades:

\begin{equation}
    \frac{\exp(V_{ni} + c X_{kn})}
    {\sum_{j \in C_n} \exp(V_{nj} + c X_{kn})}
    =
    \frac{\exp(V_{ni})}
    {\sum_{j \in C_n} \exp(V_{nj})}.
\end{equation}

Este punto abre una decisión metodológica que todavía no se encuentra cerrada. En particular, queda por definir cuál es la forma más clara y correcta de reportar los coeficientes asociados a variables comunes: tratarlas directamente como efectos relativos respecto de una alternativa base, o imponer una restricción que permita presentar coeficientes reportables para las tres alternativas. Por esta razón, en este reporte se presentan dos rutas posibles de especificación y reporte.

\subsection{Ruta A: \texttt{BIP} como alternativa base}

La primera ruta utiliza \texttt{BIP} como alternativa base para las variables comunes al viaje o a la zona. Bajo esta parametrización, para una variable común \(X_{kn}\), se fija:

\begin{equation}
    \beta_{k,\texttt{BIP}} = 0,
\end{equation}

y se estiman los coeficientes asociados a \texttt{QR\_RED} y \texttt{QR\_OTHER}. De forma simplificada, la contribución de una variable común a las utilidades queda expresada como:

\begin{align}
    V_{n,\texttt{BIP}}
    &=
    \cdots, \\
    V_{n,\texttt{QR\_RED}}
    &=
    \cdots
    +
    \beta_{k,\texttt{QR\_RED}} X_{kn}, \\
    V_{n,\texttt{QR\_OTHER}}
    &=
    \cdots
    +
    \beta_{k,\texttt{QR\_OTHER}} X_{kn}.
\end{align}

En esta ruta, los coeficientes de \texttt{QR\_RED} y \texttt{QR\_OTHER} se interpretan directamente como cambios en la utilidad relativa de cada alternativa QR respecto de \texttt{BIP}. Para las variables que sí varían entre alternativas, como tiempos y transbordos, se mantienen coeficientes específicos para \texttt{BIP}, \texttt{QR\_RED} y \texttt{QR\_OTHER}, ya que en ese caso el coeficiente de \texttt{BIP} es directamente identificable.

La ventaja de esta ruta es que corresponde a una normalización estándar y directa. Su limitación es que, para las variables comunes, \texttt{BIP} queda implícito como categoría de referencia, lo que puede dificultar una presentación simétrica de resultados en formato wide.

\subsection{Ruta B: coeficientes reportables para las tres alternativas}

La segunda ruta mantiene la misma información empírica, pero reparametriza las variables comunes para poder reportar coeficientes asociados a \texttt{BIP}, \texttt{QR\_RED} y \texttt{QR\_OTHER}. Para cada variable común \(X_{kn}\), se impone una restricción de suma cero:

\begin{equation}
    \beta_{k,\texttt{BIP}}
    +
    \beta_{k,\texttt{QR\_RED}}
    +
    \beta_{k,\texttt{QR\_OTHER}}
    =
    0.
\end{equation}

Bajo esta restricción, el coeficiente de \texttt{BIP} para una variable común no se estima libremente, sino que se deriva a partir de los otros dos coeficientes:

\begin{equation}
    \beta_{k,\texttt{BIP}}
    =
    -\beta_{k,\texttt{QR\_RED}}
    -\beta_{k,\texttt{QR\_OTHER}}.
\end{equation}

Esta parametrización permite presentar los resultados en formato wide, con columnas para \texttt{BIP}, \texttt{QR\_RED} y \texttt{QR\_OTHER}, manteniendo la identificación del modelo. Sin embargo, debe aclararse que, para las variables comunes, el coeficiente de \texttt{BIP} es un coeficiente derivado bajo la restricción de suma cero, no un parámetro libre adicional.

\subsection{Relación entre ambas rutas}

Las dos rutas anteriores son formas equivalentes de representar el mismo contenido empírico del modelo. En ambos casos se identifican diferencias relativas entre alternativas, y no efectos absolutos independientes para variables que son comunes dentro de una observación.

En la estimación realizada, ambas parametrizaciones entregan la misma verosimilitud final, el mismo AIC, el mismo BIC y el mismo número de parámetros estimados. Por lo tanto, la elección entre ambas no corresponde a una diferencia sustantiva en ajuste estadístico, sino a una decisión de interpretación y presentación. La Ruta A es más directa como normalización econométrica, mientras que la Ruta B facilita la presentación de resultados en formato wide con una columna para cada alternativa.

\section{Variables del modelo}

\subsection{Estructura general de las variables}

El modelo se estima a nivel de viaje individual. Cada observación corresponde a un viaje realizado en el sistema de transporte público y contiene, además de la alternativa de pago observada, un conjunto de atributos asociados al viaje, al contexto temporal y al entorno territorial de la zona de origen.

Las variables utilizadas en la especificación se pueden agrupar en tres grandes bloques. El primer bloque corresponde a atributos directamente asociados al viaje y a las alternativas disponibles, tales como tiempos de viaje, tiempos de espera y número de transbordos. Estas variables pueden tomar valores distintos para \texttt{BIP}, \texttt{QR\_RED} y \texttt{QR\_OTHER}, por lo que entran como variables alternativas-específicas.

El segundo bloque corresponde a variables temporales, de demanda y de oferta operacional. Estas variables describen el contexto en que ocurre el viaje, por ejemplo el año de observación, la franja horaria, el nivel de demanda local y la disponibilidad de infraestructura u oferta de transporte en la zona de origen. En general, estas variables son comunes a las tres alternativas dentro de una misma observación.

El tercer bloque corresponde a variables territoriales construidas a partir de fuentes externas, principalmente Censo 2024, OpenStreetMap y la Encuesta Origen Destino de Santiago 2012. Estas variables caracterizan el entorno socioeconómico, urbano, de infraestructura y de composición histórica de la zona de origen del viaje. Para integrarlas al modelo, cada viaje se asocia a una zona territorial \texttt{ZONA777}, y las variables zonales se incorporan según la zona de inicio del viaje.

En términos de escala, las variables con prefijo \texttt{LOG\_} corresponden a transformaciones logarítmicas del tipo \(\log(1+x)\). Esta transformación se utiliza para reducir la asimetría de variables de conteo o densidad operacional y para permitir una interpretación asociada a cambios proporcionales. Por otra parte, las variables con sufijo \texttt{\_z} corresponden a variables estandarizadas, es decir, transformadas restando su media y dividiendo por su desviación estándar. En estos casos, el coeficiente se interpreta como el cambio en utilidad asociado a un aumento de una desviación estándar en la variable.

Finalmente, es importante distinguir entre variables alternativas-específicas y variables comunes. Las primeras pueden tener valores distintos para cada alternativa de pago dentro de un mismo viaje, mientras que las segundas tienen el mismo valor para \texttt{BIP}, \texttt{QR\_RED} y \texttt{QR\_OTHER}. Esta distinción es relevante para la interpretación de los coeficientes y para la decisión metodológica discutida en la sección anterior.

\subsection{Variables de viaje, tiempo y transbordos}

El primer grupo de variables corresponde a atributos directamente asociados al viaje y a las alternativas de pago. Estas variables buscan capturar diferencias en el costo temporal y operacional asociado a cada alternativa disponible. En la base de estimación aparecen con sufijos por alternativa, por ejemplo \texttt{TVH\_BIP}, \texttt{TVH\_QR\_RED} y \texttt{TVH\_QR\_OTHER}.

Estas variables se construyen a partir de los viajes observados, agregando información por par origen--destino y tipo de pago. Para cada alternativa, se calcula el promedio observado de los atributos de viaje entre viajes comparables. Para la alternativa efectivamente elegida, se utiliza una corrección tipo \textit{leave-one-out}, de modo que el viaje observado no contribuya a construir su propio atributo alternativo. Esto evita que la variable de contexto de la alternativa elegida incorpore mecánicamente el mismo viaje que se está modelando.

Las variables de tiempo utilizadas son:

\begin{itemize}
    \item \texttt{T\_VEH}: tiempo en vehículo de la alternativa, construido a partir de la suma de los tiempos en vehículo de las etapas del viaje. En la base aparece como \texttt{TVH\_*} y se mide en segundos.
    \item \texttt{T\_ESPERA\_INI}: tiempo de espera inicial de la alternativa. En la base aparece como \texttt{TEI\_*} y se mide en segundos.
    \item \texttt{T\_ESPERA\_TRASB}: tiempo de espera asociado a transbordos. En la base aparece como \texttt{TET\_*} y se mide en segundos.
    \item \texttt{N\_TRASB}: número de transbordos de la alternativa. En la base aparece como \texttt{NTR\_*} y corresponde a una variable de conteo.
\end{itemize}

A diferencia de las variables territoriales o de contexto, estas variables pueden tomar valores distintos para \texttt{BIP}, \texttt{QR\_RED} y \texttt{QR\_OTHER} dentro de una misma observación. Por esta razón, sus coeficientes son directamente identificables para las tres alternativas, incluyendo \texttt{BIP}. En términos de interpretación, coeficientes negativos para tiempos de espera o transbordos indican una menor utilidad de la alternativa a medida que aumenta el costo temporal u operacional asociado a ella.

\subsection{Variables temporales, demanda y oferta operacional}

El segundo grupo de variables describe el contexto temporal, la intensidad de demanda local y la oferta operacional disponible en la zona de origen del viaje. A diferencia de las variables de tiempo y transbordos descritas en la subsección anterior, estas variables son comunes a las tres alternativas dentro de una misma observación. Es decir, para un viaje dado, el valor de estas variables es el mismo al evaluar \texttt{BIP}, \texttt{QR\_RED} y \texttt{QR\_OTHER}.

Las variables temporales permiten controlar por diferencias sistemáticas entre años y franjas horarias. La variable \texttt{ANIO\_2025} identifica los viajes realizados durante el año 2025 dentro de la muestra conjunta 2024--2025. Las variables \texttt{LAB\_PM}, \texttt{LAB\_PT} y \texttt{NO\_LAB} describen el tipo de franja temporal del viaje. En particular, \texttt{LAB\_PM} toma valor uno para viajes en día laboral durante la punta mañana, definida por las horas 6, 7 y 8; \texttt{LAB\_PT} toma valor uno para viajes en día laboral durante la punta tarde, definida por las horas 18 y 19; y \texttt{NO\_LAB} identifica viajes en días no laborales. La categoría omitida corresponde a \texttt{LAB\_VALLE}, es decir, viajes en día laboral fuera de las horas punta definidas.

La variable de demanda local corresponde a \texttt{LOG\_DEMAND}. Esta variable busca capturar la intensidad de viajes observada en el entorno temporal y espacial del viaje. Para construirla, los viajes se agrupan por año de observación, zona de inicio y franja horaria. Luego se cuenta el número de viajes observados en cada celda año--zona--franja. Para evitar que el viaje observado contribuya mecánicamente a su propio contexto de demanda, se utiliza una corrección tipo \textit{leave-one-out}, restando una unidad al conteo de la celda y truncando el resultado en cero. Finalmente, la variable incluida en el modelo corresponde a:

\begin{equation}
    \texttt{LOG\_DEMAND}
    =
    \log(1 + \texttt{N\_VIAJES\_ZONA\_INICIO\_FRANJA}).
\end{equation}

Por lo tanto, esta variable no representa la demanda total del sistema, sino una proxy de intensidad local de demanda en la zona de origen y franja temporal del viaje, calculada separadamente para cada año de observación.

Las variables de oferta operacional describen la disponibilidad de infraestructura o servicios de transporte en la zona de origen. La variable \texttt{LOG\_BUS\_STOP\_DENSITY} se construye a partir del número de paraderos de bus únicos asociados a cada zona. Estos paraderos se asignan a zonas mediante un cruce paradero--zona, y luego se calcula la densidad dividiendo el número de paraderos por el área de la zona en kilómetros cuadrados. La variable final corresponde a \(\log(1+\texttt{BUS\_STOP\_DENSITY})\).

La variable \texttt{LOG\_BUS\_LINE\_COUNT} busca capturar la variedad de servicios de bus disponibles en la zona y franja horaria. Para construirla se utilizan perfiles de oferta de bus con información de \texttt{ServicioUsuario}, \texttt{Paradero} y \texttt{timestamp}. La variable \texttt{ServicioUsuario} se usa como proxy de servicio o línea de bus. Primero, los registros se agregan a nivel horario y se asocian los paraderos a zonas mediante el cruce paradero--zona. Luego, para cada combinación de zona, franja y hora, se cuenta el número de servicios únicos observados. Finalmente, se promedia este conteo horario dentro de cada franja, obteniendo \texttt{BUS\_LINE\_COUNT}. La variable incluida en el modelo corresponde a \(\log(1+\texttt{BUS\_LINE\_COUNT})\).

La variable \texttt{LOG\_METRO\_LINE\_COUNT} representa la disponibilidad de líneas de Metro en la zona y franja horaria. Se construye a partir de información GTFS de Metro/Red, incluyendo calendario, frecuencias, servicios activos y estaciones. Para cada fecha y hora se identifican los servicios activos; luego se cruza la cobertura estación--línea con la asignación estación--zona. Para cada combinación de zona, franja y hora se cuenta el número de líneas de Metro activas asociadas a estaciones de la zona, y posteriormente se promedia este conteo dentro de la franja. La variable final corresponde a \(\log(1+\texttt{METRO\_LINE\_COUNT})\).

En conjunto, estas variables permiten controlar por diferencias de contexto temporal, intensidad de demanda y disponibilidad de oferta de transporte que podrían estar asociadas a la elección del medio de pago. Dado que son variables comunes a las alternativas dentro de una misma observación, su interpretación depende de la parametrización adoptada para variables comunes, discutida en la sección anterior.

\subsection{Variables sociodemográficas del Censo 2024}

Las variables sociodemográficas provienen del Censo 2024 y se incorporan al modelo a nivel de zona de origen del viaje. Para ello se utilizaron dos rutas complementarias de procesamiento. La primera corresponde a variables derivadas desde la base agregada de manzana-entidad del Censo 2024, la cual sí contiene una referencia espacial fina. Estas unidades censales se cruzan espacialmente con la zonificación \texttt{ZONA777}. Cuando los límites de las unidades censales no coinciden exactamente con los límites de \texttt{ZONA777}, la asignación se realiza mediante una agregación/interpolación areal, usando la intersección espacial entre ambas geometrías. Luego, los conteos censales se agregan a nivel de \texttt{ZONA777} y las proporciones se calculan después de sumar numeradores y denominadores, evitando promediar directamente porcentajes de unidades de distinto tamaño.

La segunda ruta corresponde a variables construidas desde los microdatos comunales del Censo 2024.
Estos microdatos contienen información individual y de hogares más detallada, pero no incluyen una
localización fina a nivel de manzana o entidad, sino identificadores comunales. Por lo tanto, no
pueden ser unidos directamente a \texttt{ZONA777}. Para utilizarlos, se implementó una estrategia
operacional de espacialización intra-comunal inspirada en la discusión metodológica de
E. Graells-Garrido sobre asignación espacial de viviendas censales sin ubicación fina.\footnote{
E. Graells-Garrido, ``Cuando los datos no tienen ubicación: un método para asignar viviendas
censales'', Datagramas, 2026. Disponible en:
\path{https://datagramas.cl/2026/01/cuando-los-datos-no-tienen-ubicacion-un-metodo-para-asignar-viviendas-censales/}.
}
La idea general es combinar microdatos comunales con información agregada a unidades espaciales más
pequeñas, de modo que la distribución espacial resultante sea consistente con los agregados
observados.

En este trabajo se utiliza una versión simplificada y operacional de esa idea: se construyen perfiles
a partir de los microdatos, se asignan probabilísticamente a unidades manzana-entidad
(\texttt{MANZENT}) usando restricciones de capacidad y luego se agregan a \texttt{ZONA777} mediante
el mismo cruce espacial del pipeline censal. A diferencia del ejemplo metodológico original, aquí no
se implementa el procedimiento completo de refinamiento espacial, por lo que las variables
microespacializadas deben interpretarse como aproximaciones territoriales consistentes con los
agregados disponibles, no como ubicaciones observadas de los hogares.

El bloque censal incorporado en las especificaciones candidatas contiene seis variables. La primera es \path{share_cine18_universitaria_o_mas_micro_z}, que mide la proporción espacializada de población de 18 años o más con educación universitaria o superior. Esta variable proviene de los microdatos espacializados y actúa como proxy de capital educativo o nivel socioeconómico del entorno de origen. La segunda es \path{share_hacinamiento_z}, definida como la proporción de viviendas particulares ocupadas clasificadas como hacinadas. La tercera es \path{share_discapacidad_z}, que corresponde a la proporción de personas con discapacidad. La cuarta es \path{share_inmigrantes_z}, definida como la proporción de población inmigrante.

Adicionalmente, se incorporan dos controles sociodemográficos territoriales. La variable \path{share_mujeres_z} mide la proporción de mujeres en la zona de origen. La variable \path{share_asistencia_parv_z} mide la proporción de población en edad preescolar que asiste a educación parvularia. Esta última no debe interpretarse como una medida exacta de asistencia a sala cuna, sino como una proxy censal más amplia de asistencia a educación parvularia. Todas estas variables se estandarizan, por lo que sus coeficientes se interpretan como cambios en utilidad asociados a un aumento de una desviación estándar en la característica territorial.

Durante el proceso de especificación también se evaluaron otras variables censales que no fueron incluidas en el bloque principal. Entre ellas se encuentran:

\begin{itemize}
    \item \path{prom_edad_z}, \path{prom_escolaridad18_z}, \path{share_internet_z}, \path{share_serv_compu_z}, \path{share_serv_tel_movil_z} y \path{share_analfabet_z}.
    \item \path{share_cine18_universitaria_micro_z}, \path{share_cine18_terciaria_corta_micro_z}, \path{share_cine18_postgrado_micro_z} y \path{share_independiente_micro_z}.
\end{itemize}

Estas variables se dejaron fuera por una combinación de menor aporte incremental, señales menos robustas o significativas, redundancia con variables ya seleccionadas, o una interpretación sustantiva menos directa para el objetivo del modelo. En particular, el bloque final privilegia una especificación parsimoniosa y defendible: educación superior como proxy socioeconómica, hacinamiento como vulnerabilidad habitacional, discapacidad como dimensión de accesibilidad social, inmigración como composición territorial, y controles adicionales de composición por sexo y asistencia parvularia.

Estas variables deben interpretarse como atributos del entorno de origen del viaje, no como características individuales necesariamente observadas para cada persona que viaja. Por ejemplo, un coeficiente asociado a \texttt{share\_inmigrantes\_z} no indica que una persona inmigrante tenga mayor o menor propensión a usar cierta alternativa de pago, sino que los viajes iniciados en zonas con mayor presencia relativa de población inmigrante presentan diferencias sistemáticas en la utilidad relativa de las alternativas.

\subsection{Variables de OpenStreetMap}

El tercer grupo de variables territoriales proviene de OpenStreetMap (OSM). A diferencia de las variables censales, que describen composición sociodemográfica de la zona, las variables OSM buscan capturar atributos del entorno construido, disponibilidad de servicios urbanos y elementos de infraestructura próximos al origen del viaje. En este reporte se tratan todas estas variables como parte de un mismo bloque OSM, incluyendo tanto equipamientos urbanos como variables de microinfraestructura de transporte.

El procesamiento comienza con una extracción de objetos OSM dentro del área de estudio. Para cada objeto se utiliza su geometría representativa y se realiza un cruce espacial con la zonificación \texttt{ZONA777}. Luego, para cada zona se cuentan los objetos que cumplen con una combinación específica de llave y valor OSM, por ejemplo \texttt{amenity=school} o \texttt{railway=subway\_entrance}. Estos conteos se dividen por el área de la zona en kilómetros cuadrados, obteniendo densidades del tipo:

\begin{equation}
    \texttt{OSM\_FEATURE\_DENSITY\_KM2}
    =
    \frac{\texttt{N\_OBJETOS\_OSM\_FEATURE}}{\texttt{AREA\_ZONA\_KM2}}.
\end{equation}

Finalmente, las densidades se estandarizan, por lo que las variables incluidas en el modelo tienen sufijo \texttt{\_z}. En consecuencia, sus coeficientes se interpretan como el cambio en utilidad asociado a un aumento de una desviación estándar en la densidad del atributo OSM correspondiente.

El bloque OSM principal incluye cinco variables de equipamiento urbano. La variable \path{osm_amenity_school_density_km2_z} mide la densidad de objetos etiquetados como \texttt{amenity=school}; \path{osm_amenity_university_density_km2_z} mide la densidad de objetos \texttt{amenity=university}; \path{osm_leisure_playground_density_km2_z} mide la densidad de áreas o puntos asociados a juegos infantiles; \path{osm_leisure_sports_centre_density_km2_z} mide la densidad de centros deportivos; y \path{osm_shop_convenience_density_km2_z} mide la densidad de comercio de conveniencia. En conjunto, estas variables aproximan centralidad barrial, intensidad de servicios locales, usos del suelo y composición funcional del entorno urbano.

Además del bloque principal, el modelo seleccionado incorpora dos variables OSM de infraestructura de acceso al transporte. La primera es \path{osm_transport_shelter_yes_density_km2_z}. Esta variable no corresponde a todos los objetos OSM con \texttt{shelter=yes}, sino a objetos filtrados como relevantes para transporte, por ejemplo paraderos, plataformas, estaciones o elementos con etiquetas asociadas a transporte público. Por lo tanto, se interpreta como una proxy de presencia de refugio o techo en puntos de espera de transporte, no como un inventario oficial completo de paraderos con refugio. La segunda es \path{osm_railway_subway_entrance_density_km2_z}, construida a partir de objetos \texttt{railway=subway\_entrance}. Esta variable aproxima la presencia física de accesos a Metro en la zona, y no simplemente la existencia de una estación o línea de Metro.

Durante el proceso de exploración también se evaluaron otras variables OSM que no fueron incorporadas en la especificación principal. Entre ellas se encuentran \path{osm_leisure_park_density_km2_z}, \path{osm_shop_supermarket_density_km2_z}, \path{osm_amenity_restaurant_density_km2_z}, \path{osm_shop_mall_density_km2_z}, \path{osm_amenity_pharmacy_density_km2_z}, \path{osm_office_company_density_km2_z}, \path{osm_office_government_density_km2_z} y \path{osm_transport_bench_yes_density_km2_z}. También se revisaron variables de conectividad peatonal o vial, como cruces y semáforos, pero se consideraron demasiado transversales para el objetivo principal del modelo. Estas variables se dejaron fuera por una combinación de menor aporte incremental, posible redundancia con variables de oferta ya incluidas, menor estabilidad estadística o interpretación menos directa respecto de la elección del medio de pago.

La interpretación de las variables OSM requiere una cautela adicional. OpenStreetMap es una fuente colaborativa y su cobertura puede variar espacialmente según el nivel de mapeo de cada zona. Por ello, estas variables no deben entenderse como mediciones administrativas exhaustivas, sino como proxies observables del entorno construido y de la infraestructura mapeada. Aun así, su valor para este modelo está en capturar diferencias territoriales que no aparecen en las variables operacionales tradicionales, especialmente en dimensiones de equipamiento urbano y experiencia de acceso al transporte.

\subsection{Variables de la Encuesta Origen Destino 2012}

Además de las fuentes contemporáneas, se evalúan variables provenientes de la Encuesta Origen Destino de Santiago 2012 (EOD 2012). Esta encuesta contiene información de hogares, personas y viajes, incluyendo variables de ingreso de hogar, tenencia de vehículos y patrones de movilidad. En este reporte la EOD no se utiliza como medición contemporánea directa del período 2024--2025, sino como una fuente histórica para construir proxies territoriales relativamente persistentes.

Para compatibilizar la EOD 2012 con la base de modelación, las variables se agregan primero a la zonificación original de la encuesta y luego se transfieren a \texttt{ZONA777} mediante un cruce espacial por intersección de áreas. En los casos en que una zona \texttt{ZONA777} intersecta varias zonas EOD, las variables se calculan como promedios ponderados por la proporción de área intersectada. De este modo, cada viaje recibe el valor EOD asociado a su zona de origen \texttt{ZONA777}.

La dimensión EOD evaluada en las especificaciones candidatas corresponde a composición histórica de ingreso del hogar. Para construirla, los hogares de la EOD 2012 se clasifican según percentiles ponderados de ingreso, usando el factor de expansión muestral de la encuesta. Este factor indica cuántos hogares de la población representa cada hogar encuestado, por lo que permite calcular distribuciones más consistentes con el diseño muestral de la EOD. A partir de esta clasificación se definen dos grupos operativos: \texttt{ABC1\_proxy}, correspondiente al tramo superior de ingreso, y \texttt{D+E\_proxy}, correspondiente a los tramos bajos. Estas etiquetas no corresponden a una clasificación socioeconómica oficial, sino a proxies históricos construidos exclusivamente a partir del ingreso declarado en la EOD 2012.

Se reportan dos formas alternativas de incorporar esta información. La primera utiliza variables continuas: \path{eod2012_share_hogares_abc1_income_proxy_z} y \path{eod2012_share_hogares_de_income_proxy_z}, que miden la proporción zonal de hogares en el tramo superior y en los tramos bajos de ingreso, respectivamente. En términos simples, estas variables mantienen la intensidad de la composición de ingreso: una zona puede tener 5\%, 20\% o 60\% de hogares en un determinado tramo, y esa diferencia gradual se conserva en el modelo.

La segunda forma utiliza variables binarias de alta concentración: \path{eod2012_dummy_high_abc1_income_proxy_z} y \path{eod2012_dummy_high_de_income_proxy_z}. Estas variables no distinguen todos los niveles posibles de la proporción zonal, sino que separan las zonas con concentración especialmente alta del resto. En particular, toman valor uno cuando la zona se ubica en el 25\% superior de la distribución zonal de la proporción correspondiente, y cero en caso contrario. Por lo tanto, la versión continua captura variación gradual en composición de ingreso, mientras que la versión dummy captura pertenencia a un grupo de zonas de alta concentración relativa.

Estas variables deben interpretarse con cautela. Primero, provienen de una encuesta de 2012, por lo que no capturan cambios recientes en composición socioeconómica, patrones residenciales o adopción tecnológica. Segundo, se incorporan como atributos agregados de la zona de origen, no como información individual del usuario que realiza el viaje. Tercero, al estar construidas a partir de ingreso histórico, pueden estar correlacionadas con otras dimensiones territoriales, como educación, hacinamiento, motorización o centralidad urbana. Su rol en el reporte es evaluar si una señal histórica de composición socioeconómica territorial agrega información interpretable al modelo, no reemplazar mediciones contemporáneas del Censo 2024.

\subsection{Resumen de variables utilizadas}

La Tabla~\ref{tab:variables_resumen} resume las variables incluidas en las especificaciones candidatas reportadas. La columna de tratamiento indica si la variable es alternativa-específica o común al viaje, y resume la transformación aplicada antes de la estimación.

\begin{landscape}
\begingroup
\scriptsize
\setlength{\tabcolsep}{3pt}
\begin{longtable}{p{0.16\linewidth} p{0.10\linewidth} p{0.35\linewidth} p{0.18\linewidth} p{0.16\linewidth}}
\caption{Resumen de variables utilizadas en el modelo}
\label{tab:variables_resumen}\\
\toprule
Variable & Bloque & Definición & Fuente & Tratamiento \\
\midrule
\endfirsthead

\toprule
Variable & Bloque & Definición & Fuente & Tratamiento \\
\midrule
\endhead

\bottomrule
\endfoot

\texttt{T\_VEH} & Viaje & Tiempo en vehículo de la alternativa. & Viajes observados & Alt.-específica; segundos \\
\texttt{T\_ESPERA\_INI} & Viaje & Tiempo de espera inicial de la alternativa. & Viajes observados & Alt.-específica; segundos \\
\texttt{T\_ESPERA\_TRASB} & Viaje & Tiempo de espera asociado a transbordos. & Viajes observados & Alt.-específica; segundos \\
\texttt{N\_TRASB} & Viaje & Número de transbordos de la alternativa. & Viajes observados & Alt.-específica; conteo \\

\texttt{ANIO\_2025} & Temporal & Indicador de viaje observado en 2025. & Base de viajes & Común; dummy \\
\texttt{LAB\_PM} & Temporal & Día laboral en punta mañana. & Base de viajes & Común; dummy \\
\texttt{LAB\_PT} & Temporal & Día laboral en punta tarde. & Base de viajes & Común; dummy \\
\texttt{NO\_LAB} & Temporal & Día no laboral. & Base de viajes & Común; dummy \\

\texttt{LOG\_DEMAND} & Contexto & Intensidad local de viajes por año, zona de origen y franja. & Viajes observados & Común; log, leave-one-out \\
\texttt{BUS stops} & Contexto & Densidad de paraderos de bus en la zona de origen. & Paraderos y \texttt{ZONA777} & Común; log \\
\texttt{Bus lines} & Contexto & Número promedio de servicios de bus disponibles por zona y franja. & Oferta bus & Común; log \\
\texttt{Metro lines} & Contexto & Número promedio de líneas de Metro activas por zona y franja. & GTFS Metro/Red & Común; log \\

Educación univ. 18+ & Censo & Proporción espacializada de población 18+ con educación universitaria o superior. & Microdatos Censo 2024 espacializados & Común; estandarizada \\
Hacinamiento & Censo & Proporción de viviendas particulares ocupadas con hacinamiento. & Censo 2024 manzana-entidad & Común; estandarizada \\
Discapacidad & Censo & Proporción de personas con discapacidad. & Censo 2024 manzana-entidad & Común; estandarizada \\
Inmigrantes & Censo & Proporción de población inmigrante. & Censo 2024 manzana-entidad & Común; estandarizada \\
Mujeres & Censo & Proporción de mujeres en la zona de origen. & Censo 2024 manzana-entidad & Común; estandarizada \\
Asistencia parvularia & Censo & Proporción de población en edad preescolar que asiste a educación parvularia. & Censo 2024 manzana-entidad & Común; estandarizada \\

Escuelas OSM & OSM & Densidad de objetos \texttt{amenity=school}. & OpenStreetMap & Común; densidad/km\(^2\), est. \\
Universidades OSM & OSM & Densidad de objetos \texttt{amenity=university}. & OpenStreetMap & Común; densidad/km\(^2\), est. \\
Playgrounds OSM & OSM & Densidad de juegos infantiles. & OpenStreetMap & Común; densidad/km\(^2\), est. \\
Centros deportivos OSM & OSM & Densidad de centros deportivos. & OpenStreetMap & Común; densidad/km\(^2\), est. \\
Convenience OSM & OSM & Densidad de comercio de conveniencia. & OpenStreetMap & Común; densidad/km\(^2\), est. \\
Refugios transporte OSM & OSM & Densidad de objetos de transporte con \texttt{shelter=yes}. & OpenStreetMap & Común; densidad/km\(^2\), est. \\
Accesos Metro OSM & OSM & Densidad de accesos físicos a Metro. & OpenStreetMap & Común; densidad/km\(^2\), est. \\

\texttt{ABC1\_proxy} EOD & EOD 2012 & Proporción histórica de hogares en el tramo superior de ingreso EOD. & EOD Santiago 2012 & Común; estandarizada \\
\texttt{D+E\_proxy} EOD & EOD 2012 & Proporción histórica de hogares en tramos bajos de ingreso EOD. & EOD Santiago 2012 & Común; estandarizada \\
Alta concentración \texttt{ABC1\_proxy} & EOD 2012 & Indicador de zona en el 25\% superior de concentración del tramo alto de ingreso. & EOD Santiago 2012 & Común; dummy est. \\
Alta concentración \texttt{D+E\_proxy} & EOD 2012 & Indicador de zona en el 25\% superior de concentración de tramos bajos de ingreso. & EOD Santiago 2012 & Común; dummy est. \\

\end{longtable}
\endgroup
\end{landscape}

\noindent
La categoría temporal omitida es \texttt{LAB\_VALLE}. Para las variables comunes al viaje, \texttt{BIP} se utiliza como alternativa de referencia; por lo tanto, los coeficientes reportados para \texttt{QR\_OTHER} y \texttt{QR\_RED} se interpretan respecto de \texttt{BIP}.

\section{Estrategia de estimación y muestra}

\subsection{Muestra de estimación}

La estimación se realiza sobre una muestra conjunta de viajes observados en los años 2024 y 2025. La unidad de observación es un viaje individual, para el cual se observa la alternativa de pago utilizada y se construyen atributos asociados a las alternativas disponibles. El conjunto de elección considerado en el modelo está compuesto por tres alternativas: \texttt{BIP}, \texttt{QR\_RED} y \texttt{QR\_OTHER}.

La muestra utilizada en la estimación principal corresponde a una submuestra estratificada del 2\% de la base interanual enriquecida. El submuestreo se realiza preservando la composición por partición temporal y por alternativa de pago observada. Es decir, dentro de cada combinación de período de datos y alternativa elegida se selecciona aleatoriamente la misma fracción de viajes, usando una semilla fija para asegurar reproducibilidad. De esta forma, la muestra conserva la presencia relativa de los distintos períodos y de las tres alternativas de pago: \texttt{BIP}, \texttt{QR\_RED} y \texttt{QR\_OTHER}.

Esta estrategia es relevante porque las alternativas de pago no tienen necesariamente la misma frecuencia en los datos. Un muestreo aleatorio simple podría alterar la proporción relativa de \texttt{BIP}, \texttt{QR\_RED} y \texttt{QR\_OTHER}, afectando la estabilidad de la estimación. Además, para evaluar que los resultados no dependan excesivamente del tamaño de muestra, se estimaron sensibilidades con muestras mayores en Larch. A diferencia de Biogeme en esta especificación, Larch permitió correr el mismo tipo de modelo con una mayor cantidad de datos sin generar problemas de memoria, por lo que se estimaron versiones al 5\% y al 10\% de la base. Estas corridas mostraron métricas de ajuste y patrones generales consistentes con la muestra de 2\%, lo que respalda el uso de la muestra reducida como base principal para la estimación y reporte en Biogeme.

En la corrida final reportada, la muestra contiene \(466{,}639\) observaciones y no se excluyen observaciones adicionales durante la estimación. El uso de una muestra interanual permite comparar patrones de elección entre 2024 y 2025, controlando explícitamente por el año de observación mediante la variable \texttt{ANIO\_2025}. De esta forma, el modelo no se estima separadamente por año, sino sobre una base pooled que permite capturar diferencias temporales sistemáticas dentro de una misma especificación.

La alternativa \texttt{BIP} se utiliza como referencia para las constantes específicas de alternativa. Por lo tanto, las constantes estimadas para \texttt{QR\_RED} y \texttt{QR\_OTHER} representan diferencias promedio de utilidad respecto de \texttt{BIP}, luego de controlar por el resto de variables incluidas en el modelo.

\subsection{Construcción de la base enriquecida}

La base de estimación se construye a partir de los viajes observados y se enriquece con información operacional, temporal y territorial. Primero, para cada viaje se incorporan variables alternativas-específicas de tiempo en vehículo, tiempo de espera inicial, tiempo de espera en transbordos y número de transbordos. Estas variables se calculan para cada alternativa de pago usando viajes comparables por origen, destino y tipo de pago, aplicando una corrección tipo \textit{leave-one-out} para evitar que el viaje observado contribuya mecánicamente a sus propios atributos.

En segundo lugar, se agregan variables de contexto temporal, demanda y oferta operacional. Las variables temporales identifican el año de observación y la franja horaria del viaje. La demanda local se calcula a partir del número de viajes observados en la zona de origen y franja temporal correspondiente. Las variables de oferta de transporte resumen la disponibilidad de paraderos de bus, servicios de bus y líneas de Metro en la zona de origen, combinando información espacial de paraderos, perfiles de oferta de bus y datos GTFS de Metro/Red.

En tercer lugar, cada viaje se asocia a una zona de origen \texttt{ZONA777}. Esta zona funciona como unidad espacial común para unir atributos territoriales provenientes de fuentes externas. El enriquecimiento territorial se realiza usando la zona donde comienza el viaje, no la zona de destino. Por lo tanto, las variables censales y OSM incluidas en el modelo describen el entorno residencial, urbano y de acceso al sistema de transporte en el punto de inicio del viaje.

Finalmente, se incorporan variables sociodemográficas del Censo 2024 y variables de OpenStreetMap. Las variables censales provienen de dos rutas: agregación espacial directa desde la base manzana-entidad hacia \texttt{ZONA777}, y espacialización de microdatos comunales hacia unidades \texttt{MANZENT} antes de agregarlos a \texttt{ZONA777}. Las variables OSM se construyen como densidades de objetos por kilómetro cuadrado dentro de cada zona. Todas las variables territoriales seleccionadas se estandarizan antes de ser incluidas en el modelo, lo que permite interpretar sus coeficientes como cambios asociados a un aumento de una desviación estándar en la característica territorial.

\subsection{Especificaciones estimadas}

La especificación central del reporte corresponde al modelo principal territorial. Este modelo combina atributos de viaje, controles temporales, variables de demanda y oferta operacional, variables sociodemográficas del Censo 2024 y variables de OpenStreetMap asociadas al entorno construido y a infraestructura de acceso al transporte. La motivación de esta especificación es evaluar si, además de los atributos operacionales del viaje, existen diferencias territoriales sistemáticas asociadas a la elección del medio de pago.

El modelo principal utiliza \texttt{BIP} como alternativa de referencia para las variables comunes al viaje o a la zona. Por lo tanto, para variables como año, franja horaria, demanda, oferta, Censo y OSM, los coeficientes de \texttt{QR\_RED} y \texttt{QR\_OTHER} se interpretan respecto de \texttt{BIP}. Esta normalización es directa desde el punto de vista de identificación del modelo y coincide con la forma usual de reportar un MNL con una alternativa base.

Además del modelo principal territorial, se estimaron sensibilidades para evaluar la robustez de la especificación. En particular, se compararon variantes con distintas combinaciones de variables OSM de infraestructura de acceso, incluyendo una especificación con accesos a Metro y otra con refugios de transporte más accesos a Metro. También se estimaron corridas en Larch con muestras de mayor tamaño para verificar que los patrones generales no dependieran exclusivamente de la muestra de 2\% usada en Biogeme.

\subsection{Criterios de evaluación}

Las especificaciones se evalúan combinando criterios estadísticos, computacionales e interpretativos. En primer lugar, se revisa la convergencia numérica del modelo. Una estimación se considera utilizable si el proceso de optimización converge sin errores, si el gradiente final es razonablemente bajo y si no aparecen problemas evidentes de identificación o parámetros inestables. Este criterio es especialmente relevante porque la especificación incluye un número elevado de parámetros y variables territoriales comunes a las alternativas.

En segundo lugar, se comparan medidas de ajuste global. La principal medida de ajuste es la log-verosimilitud final, donde valores menos negativos indican un mejor ajuste a los datos. También se reportan el criterio de información de Akaike (AIC) y el criterio de información bayesiano (BIC). Ambos penalizan la complejidad del modelo, pero BIC aplica una penalización más fuerte al número de parámetros, especialmente en muestras grandes. Por esta razón, una mejora en log-verosimilitud y AIC acompañada de un aumento leve en BIC no se interpreta automáticamente como evidencia en contra del modelo, sino como una señal de que el aporte incremental debe justificarse también por interpretabilidad y valor sustantivo.

En tercer lugar, se revisa la significancia estadística y estabilidad de los parámetros. Para los resultados principales se reportan errores estándar robustos, estadísticos \(t\) y valores \(p\). La interpretación se concentra no solo en si un parámetro es significativo, sino también en si su signo es estable, si su magnitud es razonable y si la lectura sustantiva es coherente con el fenómeno estudiado.

Finalmente, la selección del modelo principal no se basa únicamente en maximizar ajuste estadístico. Dado que el objetivo del reporte es describir un modelo interpretable de elección de medio de pago, se privilegia una especificación parsimoniosa, con variables territorialmente defendibles y con lectura útil para política pública. En particular, las variables de infraestructura de acceso se mantienen cuando aportan información sustantiva adicional sobre la experiencia de acceso al sistema, aun cuando su aporte en BIC sea más exigente de justificar.

\section{Ajuste y convergencia del modelo}

\subsection{Resultados globales de estimación}

La Tabla~\ref{tab:ajuste_convergencia} resume las métricas globales del modelo principal
territorial bajo las dos formas de parametrizar las variables comunes. La Opción A mantiene a
\texttt{BIP} como alternativa de referencia para las variables comunes, mientras que la Opción B
usa una restricción de suma cero para reportar coeficientes para las tres alternativas. Ambas
especificaciones representan el mismo contenido conductual y, por lo tanto, entregan el mismo
ajuste global.

\begin{table}[htbp]
\centering
\small
\caption{Ajuste y convergencia del modelo principal territorial}
\label{tab:ajuste_convergencia}
\begin{tabular}{lrr}
\toprule
Métrica & Opción A & Opción B \\
\midrule
Parámetros estimados & 52 & 52 \\
Observaciones & 466{,}639 & 466{,}639 \\
Observaciones excluidas & 0 & 0 \\
Log-verosimilitud inicial & -512{,}655.3 & -512{,}655.3 \\
Log-verosimilitud final & -232{,}257.1 & -232{,}257.1 \\
Test de razón de verosimilitud & 560{,}796.5 & 560{,}796.5 \\
Rho-cuadrado & 0.5470 & 0.5470 \\
Rho-cuadrado ajustado & 0.5470 & 0.5470 \\
AIC & 464{,}618.2 & 464{,}618.2 \\
BIC & 465{,}193.0 & 465{,}193.0 \\
Norma final del gradiente & 4.7659 & 4.8173 \\
Convergencia & Sí & Sí \\
\bottomrule
\end{tabular}
\end{table}

En ambos casos, el modelo converge correctamente y no excluye observaciones de la muestra de
estimación. La igualdad en log-verosimilitud, AIC y BIC confirma que las dos opciones no
corresponden a modelos sustantivamente distintos, sino a dos normalizaciones alternativas para
presentar los coeficientes de variables comunes. La diferencia principal está en la interpretación
de los parámetros: en la Opción A los coeficientes de variables comunes se leen respecto de
\texttt{BIP}, mientras que en la Opción B los coeficientes se leen como desviaciones respecto del
promedio de alternativas.

\subsection{Lectura de las métricas de ajuste}

La log-verosimilitud final mejora sustancialmente respecto de la log-verosimilitud inicial, pasando
de $-512{,}655.3$ a $-232{,}257.1$. Esto indica que las variables incluidas en la especificación
aportan capacidad explicativa relevante respecto de un modelo sin parámetros informativos. En la
misma línea, el test de razón de verosimilitud alcanza un valor de $560{,}796.5$, lo que refleja una
mejora global importante frente al modelo inicial.

El valor de rho-cuadrado es $0.5470$, lo que sugiere un ajuste alto para un modelo de elección
discreta aplicado a viajes individuales. Esta métrica no debe interpretarse como un $R^2$ de
regresión lineal, pero sí permite evaluar cuánto mejora el modelo respecto de una especificación
base. En este caso, el resultado indica que el modelo logra capturar una parte sustantiva de la
estructura observada de elección entre \texttt{BIP}, \texttt{QR\_RED} y \texttt{QR\_OTHER}.

Los criterios AIC y BIC son iguales entre ambas opciones, porque ambas tienen la misma
log-verosimilitud y el mismo número de parámetros estimados. Por lo tanto, no sirven para escoger
entre la Opción A y la Opción B. La elección entre ambas debe entenderse como una decisión
metodológica de interpretación y presentación de resultados, no como una comparación de ajuste
predictivo o estadístico.

Finalmente, ambas especificaciones reportan convergencia numérica. La Opción B requiere más
iteraciones que la Opción A debido a la parametrización con restricción de suma cero, pero alcanza
el mismo óptimo de log-verosimilitud. Por esta razón, la decisión entre ambas rutas se deja
planteada como una decisión de reporte: mantener una lectura tradicional respecto de \texttt{BIP} o
presentar coeficientes comparables para las tres alternativas.

\section{Resultados de parámetros en formato wide}

\subsection{Parámetros estimados}

Las Tablas~\ref{tab:parametros_modelo_continuo} y~\ref{tab:parametros_modelo_dummy} presentan los
parámetros de las dos especificaciones candidatas. En ambas, \texttt{BIP} se mantiene como
alternativa de referencia para las variables comunes al conjunto de alternativas. Por lo tanto, para
variables como año, franja horaria, demanda, oferta, Censo, OSM y EOD, el coeficiente de
\texttt{BIP} se fija en cero y los coeficientes de \texttt{QR\_OTHER} y \texttt{QR\_RED} se
interpretan respecto de esa alternativa base.

Esta normalización no implica que dichas variables no estén asociadas a la elección de
\texttt{BIP}. Más bien, implica que el efecto se expresa en términos relativos: un coeficiente
positivo para \texttt{QR\_RED}, por ejemplo, indica que un aumento en la variable incrementa la
utilidad de \texttt{QR\_RED} respecto de \texttt{BIP}, manteniendo todo lo demás constante. Para las
variables que sí varían entre alternativas dentro de un mismo viaje, como tiempos y transbordos,
también se estima directamente un coeficiente para \texttt{BIP}.

Se presentan los coeficientes estimados y sus estadísticos \emph{t-ratio} robustos. Una versión
extendida de las tablas, que incluye errores estándar robustos, se deja en el
Anexo~\ref{sec:anexo_parametros_completos}. Los ceros reportados para \texttt{BIP} en variables
comunes corresponden a restricciones de normalización y no a parámetros estimados.

Además, para facilitar la interpretación de magnitudes, el Anexo~\ref{sec:anexo_odds_ratios}
presenta los \emph{odds ratios} asociados a ambas especificaciones. Estos se calculan como
\(\exp(\beta)\) o, en el caso de variables de tiempo medidas en segundos, como \(\exp(60\beta)\)
para reportar el cambio asociado a un minuto adicional.

\subsubsection{Modelo con proxy EOD continuo}

La Tabla~\ref{tab:parametros_modelo_continuo} presenta la especificación que incorpora la composición
histórica de ingreso EOD mediante proporciones continuas de hogares en los tramos
\texttt{ABC1\_proxy} y \texttt{D+E\_proxy}.

\begin{landscape}
\small
\setlength{\tabcolsep}{3pt}
\begin{longtable}{p{0.27\linewidth}rrrrrr}
\caption{Parámetros del modelo con proxy EOD continuo}
\label{tab:parametros_modelo_continuo}\\
\toprule
Variable &
\texttt{BIP} Est. & \texttt{BIP} t &
\texttt{QR\_OTHER} Est. & \texttt{QR\_OTHER} t &
\texttt{QR\_RED} Est. & \texttt{QR\_RED} t \\
\midrule
\endfirsthead

\toprule
Variable &
\texttt{BIP} Est. & \texttt{BIP} t &
\texttt{QR\_OTHER} Est. & \texttt{QR\_OTHER} t &
\texttt{QR\_RED} Est. & \texttt{QR\_RED} t \\
\midrule
\endhead

\bottomrule
\endfoot

\texttt{ASC} & -- & -- & -1.9172 & -23.32 & -3.5839 & -20.81 \\
\texttt{ANIO\_2025} & 0 & -- & 0.2965 & 33.32 & 0.2332 & 12.72 \\
\texttt{LAB\_PM} & 0 & -- & 0.0976 & 6.76 & 0.4604 & 15.25 \\
\texttt{LAB\_PT} & 0 & -- & 0.1571 & 8.87 & 0.2396 & 6.89 \\
\texttt{NO\_LAB} & 0 & -- & 0.1624 & 9.23 & 0.0510 & 1.35 \\
\texttt{LOG\_DEMAND} & 0 & -- & -0.0344 & -4.31 & -0.0671 & -4.24 \\
\texttt{LOG\_BUS\_STOP\_DENSITY} & 0 & -- & 0.0008 & 0.06 & -0.0910 & -2.92 \\
\texttt{LOG\_BUS\_LINE\_COUNT} & 0 & -- & 0.0164 & 1.13 & 0.0901 & 2.91 \\
\texttt{LOG\_METRO\_LINE\_COUNT} & 0 & -- & 0.0488 & 2.39 & 0.1235 & 3.14 \\
\texttt{T\_VEH} & 0.000020 & 0.66 & 0.000071 & 2.39 & -0.000006 & -0.19 \\
\texttt{T\_ESPERA\_INI} & -0.0017 & -29.21 & -0.0017 & -29.39 & -0.0014 & -20.08 \\
\texttt{T\_ESPERA\_TRASB} & -0.0003 & -3.59 & -0.0004 & -5.60 & 0.000045 & 0.60 \\
\texttt{N\_TRASB} & -0.0597 & -1.50 & -0.1248 & -3.15 & -0.1556 & -3.83 \\
\texttt{Educación univ. 18+} & 0 & -- & 0.0815 & 8.93 & 0.3073 & 16.47 \\
\texttt{Hacinamiento} & 0 & -- & -0.0652 & -5.66 & -0.0492 & -1.79 \\
\texttt{Discapacidad} & 0 & -- & 0.0028 & 0.24 & -0.0988 & -3.67 \\
\texttt{Inmigrantes} & 0 & -- & 0.1007 & 9.91 & 0.0865 & 3.68 \\
\texttt{Mujeres} & 0 & -- & -0.0002 & -0.06 & 0.0212 & 2.37 \\
\texttt{Asistencia parvularia} & 0 & -- & 0.0053 & 0.76 & 0.0487 & 3.11 \\
\texttt{Centros deportivos OSM} & 0 & -- & -0.0092 & -2.10 & 0.0078 & 0.89 \\
\texttt{Convenience OSM} & 0 & -- & 0.0135 & 3.07 & 0.0414 & 4.69 \\
\texttt{Playgrounds OSM} & 0 & -- & 0.0283 & 5.65 & 0.0662 & 6.44 \\
\texttt{Escuelas OSM} & 0 & -- & -0.0001 & -0.03 & -0.0284 & -2.54 \\
\texttt{Universidades OSM} & 0 & -- & -0.0197 & -9.15 & -0.0244 & -5.35 \\
\texttt{Refugios transporte OSM} & 0 & -- & -0.0306 & -4.65 & 0.0375 & 2.62 \\
\texttt{Accesos Metro OSM} & 0 & -- & -0.0096 & -2.83 & -0.0370 & -5.27 \\
\texttt{ABC1\_proxy EOD} & 0 & -- & -0.0154 & -1.87 & -0.0053 & -0.34 \\
\texttt{D+E\_proxy EOD} & 0 & -- & 0.0369 & 4.53 & -0.0123 & -0.67 \\

\end{longtable}
\end{landscape}

\subsubsection{Modelo con proxy EOD dummy}

La Tabla~\ref{tab:parametros_modelo_dummy} presenta la especificación que incorpora la composición
histórica de ingreso EOD mediante indicadores de alta concentración de hogares en los tramos
\texttt{ABC1\_proxy} y \texttt{D+E\_proxy}.

\begin{landscape}
\small
\setlength{\tabcolsep}{3pt}
\begin{longtable}{p{0.27\linewidth}rrrrrr}
\caption{Parámetros del modelo con proxy EOD dummy}
\label{tab:parametros_modelo_dummy}\\
\toprule
Variable &
\texttt{BIP} Est. & \texttt{BIP} t &
\texttt{QR\_OTHER} Est. & \texttt{QR\_OTHER} t &
\texttt{QR\_RED} Est. & \texttt{QR\_RED} t \\
\midrule
\endfirsthead

\toprule
Variable &
\texttt{BIP} Est. & \texttt{BIP} t &
\texttt{QR\_OTHER} Est. & \texttt{QR\_OTHER} t &
\texttt{QR\_RED} Est. & \texttt{QR\_RED} t \\
\midrule
\endhead

\bottomrule
\endfoot

\texttt{ASC} & -- & -- & -1.9081 & -23.16 & -3.5086 & -20.47 \\
\texttt{ANIO\_2025} & 0 & -- & 0.2963 & 33.30 & 0.2329 & 12.70 \\
\texttt{LAB\_PM} & 0 & -- & 0.0955 & 6.60 & 0.4529 & 14.97 \\
\texttt{LAB\_PT} & 0 & -- & 0.1522 & 8.60 & 0.2305 & 6.62 \\
\texttt{NO\_LAB} & 0 & -- & 0.1603 & 9.11 & 0.0447 & 1.18 \\
\texttt{LOG\_DEMAND} & 0 & -- & -0.0370 & -4.65 & -0.0722 & -4.56 \\
\texttt{LOG\_BUS\_STOP\_DENSITY} & 0 & -- & -0.0022 & -0.16 & -0.1134 & -3.60 \\
\texttt{LOG\_BUS\_LINE\_COUNT} & 0 & -- & 0.0251 & 1.72 & 0.1023 & 3.29 \\
\texttt{LOG\_METRO\_LINE\_COUNT} & 0 & -- & 0.0571 & 2.80 & 0.1251 & 3.19 \\
\texttt{T\_VEH} & 0.000020 & 0.67 & 0.000069 & 2.32 & -0.000008 & -0.25 \\
\texttt{T\_ESPERA\_INI} & -0.0017 & -29.34 & -0.0017 & -29.43 & -0.0014 & -20.04 \\
\texttt{T\_ESPERA\_TRASB} & -0.0003 & -3.58 & -0.0004 & -5.58 & 0.000046 & 0.61 \\
\texttt{N\_TRASB} & -0.0644 & -1.62 & -0.1285 & -3.24 & -0.1572 & -3.86 \\
\texttt{Educación univ. 18+} & 0 & -- & 0.0784 & 8.63 & 0.3194 & 17.16 \\
\texttt{Hacinamiento} & 0 & -- & -0.0556 & -4.84 & -0.0520 & -1.88 \\
\texttt{Discapacidad} & 0 & -- & 0.0123 & 1.09 & -0.1057 & -4.02 \\
\texttt{Inmigrantes} & 0 & -- & 0.1033 & 10.24 & 0.0962 & 4.09 \\
\texttt{Mujeres} & 0 & -- & 0.0005 & 0.15 & 0.0241 & 2.68 \\
\texttt{Asistencia parvularia} & 0 & -- & 0.0094 & 1.33 & 0.0672 & 4.17 \\
\texttt{Centros deportivos OSM} & 0 & -- & -0.0086 & -1.94 & 0.0113 & 1.29 \\
\texttt{Convenience OSM} & 0 & -- & 0.0119 & 2.74 & 0.0429 & 4.94 \\
\texttt{Playgrounds OSM} & 0 & -- & 0.0283 & 5.63 & 0.0644 & 6.26 \\
\texttt{Escuelas OSM} & 0 & -- & 0.0020 & 0.39 & -0.0262 & -2.35 \\
\texttt{Universidades OSM} & 0 & -- & -0.0178 & -8.03 & -0.0211 & -4.47 \\
\texttt{Refugios transporte OSM} & 0 & -- & -0.0304 & -4.58 & 0.0328 & 2.28 \\
\texttt{Accesos Metro OSM} & 0 & -- & -0.0104 & -3.06 & -0.0340 & -4.90 \\
\texttt{Alta concentración ABC1\_proxy} & 0 & -- & -0.0264 & -4.38 & -0.0421 & -3.44 \\
\texttt{Alta concentración D+E\_proxy} & 0 & -- & -0.0026 & -0.37 & -0.0110 & -0.65 \\

\end{longtable}
\end{landscape}

\subsection{Comparación general entre especificaciones}

Los parámetros de las variables base, temporales, de viaje, Censo y OSM son estables entre ambas
especificaciones. En particular, se mantienen los signos principales asociados a año, franjas
horarias, demanda local, tiempos de espera, educación universitaria, hacinamiento, inmigración,
refugios de transporte y accesos a Metro. Esto sugiere que la forma de representar el proxy EOD no
altera de manera sustantiva la estructura central del modelo.

La diferencia principal aparece en la lectura del bloque EOD. En la especificación continua,
\texttt{D+E\_proxy} presenta una asociación positiva y significativa con \texttt{QR\_OTHER}, mientras
que \texttt{ABC1\_proxy} muestra una señal negativa solo marginal. En la especificación dummy, la
alta concentración \texttt{ABC1\_proxy} presenta coeficientes negativos y significativos tanto para
\texttt{QR\_OTHER} como para \texttt{QR\_RED}, mientras que la alta concentración
\texttt{D+E\_proxy} no presenta una señal estadísticamente robusta.

La interpretación sustantiva de estos patrones se desarrolla en la sección siguiente. En esta sección
solo se deja establecido que ambas especificaciones son comparables, convergen correctamente y
conservan patrones centrales similares, pero capturan de forma distinta la dimensión histórica de
ingreso territorial proveniente de la EOD 2012.

\section{Interpretación de signos y significancia}

\subsection{Consideraciones de interpretación}

La interpretación de los resultados se realiza usando \texttt{BIP} como alternativa de referencia
para las variables comunes al viaje o a la zona. Por lo tanto, para variables como año, franja
horaria, demanda, oferta, Censo, OSM y EOD, un coeficiente positivo de \texttt{QR\_RED} o
\texttt{QR\_OTHER} indica una mayor utilidad relativa de esa alternativa respecto de \texttt{BIP},
manteniendo constantes las demás variables del modelo. Un coeficiente negativo indica una menor
utilidad relativa respecto de \texttt{BIP}.

En el caso de variables alternativa-específicas, como tiempos de viaje, esperas y transbordos, los
coeficientes se interpretan directamente para cada alternativa. Estas variables pueden tomar valores
distintos para \texttt{BIP}, \texttt{QR\_RED} y \texttt{QR\_OTHER} dentro de una misma observación,
por lo que sus coeficientes son identificables para las tres alternativas.

Finalmente, los resultados deben interpretarse como asociaciones condicionales del modelo, no como
efectos causales. En particular, las variables territoriales pueden capturar diferencias de
composición socioespacial, intensidad de uso del sistema, disponibilidad de infraestructura y otros
factores no observados que coexisten en una misma zona. Esta cautela es especialmente relevante para
las variables EOD 2012, que se incorporan como proxies históricos de composición territorial y no
como mediciones contemporáneas directas del período 2024--2025.

\subsection{Variables de viaje}

Las variables de espera presentan los patrones más estables en ambas especificaciones candidatas. El
tiempo de espera inicial (\texttt{T\_ESPERA\_INI}) tiene coeficientes negativos, de magnitud similar
y altamente significativos para las tres alternativas. Esto es consistente con la teoría de elección
discreta: mayores tiempos de espera reducen la utilidad de la alternativa considerada.

El tiempo de espera asociado a transbordos (\texttt{T\_ESPERA\_TRASB}) también presenta coeficientes
negativos para \texttt{BIP} y \texttt{QR\_OTHER}, ambos estadísticamente significativos. En
\texttt{QR\_RED}, en cambio, el coeficiente es cercano a cero y no significativo en ambas
especificaciones. Una lectura sustantiva posible es que los viajes asociados a \texttt{QR\_RED}
presentan menor penalización por espera en transbordo, posiblemente porque el uso de la app entrega
información que permite anticipar la espera o decidir cómo utilizar ese tiempo. Esta hipótesis debe
tratarse como interpretación exploratoria y no como evidencia causal directa.

El número de transbordos (\texttt{N\_TRASB}) tiene signo negativo para las tres alternativas, aunque
su significancia es mayor para \texttt{QR\_OTHER} y \texttt{QR\_RED}. Esto indica que, manteniendo
constantes las demás variables, viajes con más transbordos tienden a reducir la utilidad asociada a
las alternativas de pago, especialmente en las alternativas QR.

Por su parte, el tiempo en vehículo (\texttt{T\_VEH}) presenta coeficientes de magnitud muy pequeña.
El único coeficiente estadísticamente significativo aparece en \texttt{QR\_OTHER}, con signo
positivo. Esta señal debe interpretarse con cautela: más que indicar una preferencia por viajes más
largos, probablemente refleja diferencias de composición entre tipos de viaje, zonas de origen,
distancias recorridas y alternativas disponibles.

\subsection{Variables temporales, demanda y oferta}

La variable \texttt{ANIO\_2025} muestra una asociación positiva con \texttt{QR\_OTHER} y
\texttt{QR\_RED} respecto de \texttt{BIP}. Esto sugiere que, en la muestra analizada, el año 2025 se
asocia con una mayor utilidad relativa de las alternativas QR. Esta señal es consistente con una
posible expansión o consolidación del uso de medios de pago QR entre 2024 y 2025.

Las variables de franja muestran diferencias relevantes entre alternativas. \texttt{LAB\_PM} se
asocia fuertemente con \texttt{QR\_RED} respecto de \texttt{BIP}, mientras que también presenta una
asociación positiva, aunque menor, con \texttt{QR\_OTHER}. \texttt{LAB\_PT} muestra una señal
similar, aunque de menor magnitud. En contraste, \texttt{NO\_LAB} se asocia positivamente con
\texttt{QR\_OTHER}, mientras que la señal para \texttt{QR\_RED} no es estadísticamente robusta.

La variable \texttt{LOG\_DEMAND} presenta coeficientes negativos y significativos para
\texttt{QR\_OTHER} y \texttt{QR\_RED} respecto de \texttt{BIP}. Esto sugiere que zonas y franjas con
mayor intensidad local de viajes tienden a estar relativamente menos asociadas a alternativas QR. Sin
embargo, esta variable debe interpretarse con cautela, porque puede reflejar tanto intensidad de uso
del sistema como composición territorial de la demanda.

En cuanto a oferta, la densidad de paraderos se asocia negativamente con \texttt{QR\_RED} y no
presenta una señal robusta para \texttt{QR\_OTHER}. El número de líneas de bus y de Metro muestra
coeficientes positivos para \texttt{QR\_RED}, y el número de líneas de Metro también es positivo para
\texttt{QR\_OTHER}. Estos resultados sugieren que la relación entre infraestructura operacional y
medio de pago no es homogénea entre alternativas QR, y probablemente captura diferencias entre tipos
de viaje, zonas de origen y estructura de la red.

\subsection{Variables censales}

La proporción de población de 18 años o más con educación universitaria o superior muestra uno de los
patrones más claros de ambos modelos. El coeficiente es positivo y altamente significativo para
\texttt{QR\_RED} respecto de \texttt{BIP}, y también positivo para \texttt{QR\_OTHER}, aunque de
menor magnitud. En términos relativos, esto indica que las zonas con mayor presencia de población
universitaria se asocian con una mayor utilidad relativa de alternativas QR, especialmente
\texttt{QR\_RED}. Este resultado es coherente con la hipótesis de que mayor nivel educacional puede
relacionarse con mayor adopción de herramientas digitales de pago o información.

El hacinamiento presenta coeficientes negativos para las alternativas QR, especialmente para
\texttt{QR\_OTHER}. En \texttt{QR\_RED} la señal también es negativa, aunque menos robusta. Este
patrón sugiere que zonas con mayor hacinamiento se asocian relativamente menos con alternativas QR,
lo que puede reflejar condiciones socioeconómicas y diferencias de acceso o adopción digital.

La proporción de personas con discapacidad muestra una asociación negativa y significativa con
\texttt{QR\_RED}, mientras que la señal para \texttt{QR\_OTHER} no es robusta en las especificaciones
actuales. En este caso, la interpretación debe ser cuidadosa, porque la variable describe el contexto
territorial de origen y no necesariamente una característica individual del usuario que realiza el
viaje.

La proporción de población inmigrante se asocia positivamente con \texttt{QR\_OTHER} y también con
\texttt{QR\_RED}. La magnitud es mayor para \texttt{QR\_OTHER}, aunque ambas señales se mantienen
positivas en los dos modelos. Esto sugiere que zonas con mayor proporción de población inmigrante
están relativamente más asociadas a alternativas QR respecto de \texttt{BIP}.

Las variables censales de composición sociodemográfica agregan dos lecturas complementarias. La proporción de mujeres en la
zona de origen no presenta una señal robusta para \texttt{QR\_OTHER}, pero sí una asociación positiva
y significativa con \texttt{QR\_RED}. Por su parte, la asistencia a educación parvularia muestra una
señal positiva y significativa para \texttt{QR\_RED}, mientras que su asociación con
\texttt{QR\_OTHER} es débil o no robusta. Dado que estas variables son territoriales, no deben
interpretarse como características individuales del usuario, sino como indicadores de composición
sociodemográfica y ciclo de vida de las zonas de origen.

\subsection{Variables OSM e infraestructura}

La densidad de universidades OSM presenta coeficientes negativos y significativos para
\texttt{QR\_OTHER} y \texttt{QR\_RED} respecto de \texttt{BIP}. Este resultado no contradice
necesariamente el efecto positivo de la educación universitaria residencial sobre \texttt{QR\_RED}.
Ambas variables capturan dimensiones distintas: la variable censal de educación universitaria
describe la composición educacional de la población residente, mientras que
\path{osm_amenity_university_density_km2_z} identifica la presencia territorial de equipamientos
universitarios. En Santiago, estos equipamientos están fuertemente concentrados en pocas comunas
centrales, por lo que esta variable también puede estar capturando centralidad urbana, concentración
de destinos educacionales y zonas con alta presencia de Metro o viajes asociados a \texttt{BIP}.

La densidad de juegos infantiles y comercio de conveniencia muestra asociaciones positivas con ambas
alternativas QR, especialmente con \texttt{QR\_RED}. Esto sugiere que ciertos entornos urbanos de
escala barrial se asocian con una mayor utilidad relativa de alternativas QR. La densidad de centros
deportivos presenta una señal menos clara: tiende a ser negativa para \texttt{QR\_OTHER} y no robusta
para \texttt{QR\_RED}.

La variable de refugios de transporte en OSM
(\path{osm_transport_shelter_yes_density_km2_z}) presenta una señal particularmente interesante:
el coeficiente es negativo y significativo para \texttt{QR\_OTHER}, pero positivo y significativo
para \texttt{QR\_RED}. Esto indica que la presencia de objetos de transporte con
\texttt{shelter=yes} no se asocia de forma homogénea con el uso de QR. Más que interpretarse como un
efecto causal de tener refugio, esta variable debe leerse como un indicador territorial de
infraestructura de parada o entorno de acceso al sistema.

La densidad de accesos a Metro
(\path{osm_railway_subway_entrance_density_km2_z}) presenta coeficientes negativos y significativos
para \texttt{QR\_OTHER} y \texttt{QR\_RED} respecto de \texttt{BIP}. Esto sugiere que zonas con
mayor presencia de accesos físicos a Metro se asocian relativamente menos con alternativas QR una vez
controladas las demás variables del modelo. La densidad de colegios también tiende a presentar una
señal negativa para \texttt{QR\_RED}, mientras que su asociación con \texttt{QR\_OTHER} no es
robusta.

\subsection{Variables EOD 2012}

Las variables de ingreso construidas desde la EOD 2012 se evalúan en dos formas alternativas. La
primera usa proporciones continuas de hogares en tramos históricos de ingreso; la segunda usa
indicadores de alta concentración zonal. Estas dos formas no responden exactamente la misma pregunta:
la especificación continua mide cambios graduales en la composición de ingreso de la zona, mientras
que la especificación dummy identifica zonas ubicadas en el cuartil superior de concentración de un
grupo específico.

En la especificación continua, la proporción histórica de hogares en los tramos bajos de ingreso
(\texttt{D+E\_proxy}) presenta una asociación positiva y significativa con \texttt{QR\_OTHER}, pero
no muestra una señal robusta para \texttt{QR\_RED}. En cambio, la proporción histórica de hogares en
el tramo superior de ingreso (\texttt{ABC1\_proxy}) presenta una señal negativa solo marginal para
\texttt{QR\_OTHER} y no robusta para \texttt{QR\_RED}. Esta lectura sugiere que el gradiente
continuo de ingreso bajo se relaciona principalmente con \texttt{QR\_OTHER}.

En la especificación dummy, la alta concentración de hogares \texttt{ABC1\_proxy} muestra
coeficientes negativos y significativos para \texttt{QR\_OTHER} y \texttt{QR\_RED}. En cambio, la
alta concentración de hogares \texttt{D+E\_proxy} no presenta una señal estadísticamente robusta.
Esto sugiere que, cuando el ingreso histórico se representa como una condición de alta concentración,
la señal más clara aparece en zonas de mayor ingreso relativo, que se asocian con menor utilidad
relativa de alternativas QR respecto de \texttt{BIP}.

Ambas lecturas son complementarias más que contradictorias. El modelo continuo permite observar una
relación gradual asociada a la presencia de hogares de ingresos bajos, especialmente para
\texttt{QR\_OTHER}; el modelo dummy, en cambio, resalta el comportamiento de zonas con concentración
alta de hogares de mayores ingresos. En cualquier caso, estas variables deben tratarse como proxies
históricos y territoriales, no como clasificación socioeconómica oficial ni como ingreso individual
del usuario.

\subsection{Síntesis}

En conjunto, los resultados muestran que la elección entre \texttt{BIP}, \texttt{QR\_OTHER} y
\texttt{QR\_RED} no solo depende de atributos operacionales del viaje, sino también de
características temporales y territoriales del origen. Las variables de espera y transbordo mantienen
los signos esperados en la mayoría de los casos, con una excepción sustantivamente interesante en
\texttt{QR\_RED}: la espera por transbordo no presenta una penalización estadísticamente robusta.

Los resultados más robustos para \texttt{QR\_RED} aparecen en la asociación positiva con educación
universitaria, puntas laborales, asistencia parvularia, proporción de mujeres y ciertos entornos OSM
de escala barrial. Para \texttt{QR\_OTHER}, destacan la asociación positiva con población inmigrante
y con el proxy continuo de hogares en tramos bajos de ingreso, junto con señales negativas asociadas
a hacinamiento, demanda local y refugios de transporte.

La comparación entre los modelos continuo y dummy muestra que la incorporación de la EOD 2012 no
altera los patrones principales del modelo, pero sí permite explorar dos formas distintas de
representar la dimensión socioeconómica histórica del territorio. Por lo tanto, la decisión entre
ambas especificaciones debería basarse principalmente en la interpretación sustantiva que se quiera
privilegiar: gradientes territoriales continuos de ingreso o identificación de zonas con alta
concentración relativa de ciertos grupos.

\section{Decisión de modelo principal y sensibilidades}

A partir de los resultados anteriores, la especificación territorial de trabajo corresponde a un
modelo MNL con \texttt{BIP} como alternativa de referencia para variables comunes, que incorpora
atributos del viaje, variables temporales, demanda y oferta operacional, variables
sociodemográficas del Censo 2024, variables OSM de entorno urbano e infraestructura de acceso, y
proxies históricos de ingreso provenientes de la EOD 2012. Esta estructura permite analizar la
elección entre \texttt{BIP}, \texttt{QR\_OTHER} y \texttt{QR\_RED} combinando atributos
operacionales del viaje con características territoriales del origen.

Dentro de esta estructura se mantienen dos especificaciones candidatas. La primera representa la
dimensión de ingreso EOD mediante proporciones continuas de hogares en los tramos
\texttt{ABC1\_proxy} y \texttt{D+E\_proxy}. La segunda utiliza indicadores de alta concentración
zonal de esos mismos grupos. Ambas especificaciones comparten el resto de variables del modelo y
mantienen patrones sustantivos similares en los bloques de viaje, temporalidad, Censo y OSM.

La decisión entre ambas representaciones no debe basarse únicamente en AIC, BIC o
log-verosimilitud. En modelos con variables territoriales correlacionadas, una variable puede aportar
valor interpretativo aun cuando la mejora incremental de ajuste sea moderada. Por esto, la elección
del modelo de reporte debe combinar tres criterios: estabilidad de signos, coherencia sustantiva e
interpretabilidad para política pública. Bajo ese criterio, la especificación continua permite leer
gradientes territoriales de composición de ingreso, mientras que la especificación dummy permite
identificar zonas con alta concentración relativa de ciertos grupos.

Respecto del bloque OSM, se mantiene la inclusión de refugios de transporte y accesos a Metro. La
variable \path{osm_transport_shelter_yes_density_km2_z} aporta una dimensión de infraestructura de
parada que no está representada por las variables operacionales tradicionales. Por su parte,
\path{osm_railway_subway_entrance_density_km2_z} agrega información sobre accesibilidad física a la
red de Metro, complementando la variable de número de líneas de Metro disponibles en la zona. Aunque
estas variables no deben interpretarse causalmente, sí entregan una lectura territorial relevante
sobre el entorno de acceso al sistema.

Respecto del bloque censal, la especificación conserva variables de composición educacional,
vulnerabilidad habitacional, discapacidad, inmigración, composición por sexo y asistencia a educación
parvularia. Estas variables no describen necesariamente al individuo que realiza el viaje, sino el
contexto territorial de la zona de origen. Su valor está en permitir que el modelo distinga patrones
de adopción de medios de pago asociados a la composición sociodemográfica del territorio.

Como sensibilidad adicional, se evaluó la variable \path{eod2012_numveh_mean_z}, que mide el número
promedio histórico de vehículos por hogar en la zona. Esta variable presentó una señal
estadísticamente informativa, especialmente para \texttt{QR\_OTHER}. Sin embargo, no se incorpora en
las especificaciones candidatas junto con los proxies de ingreso, porque representa una dimensión
conceptualmente muy cercana al nivel socioeconómico y presenta alta colinealidad empírica con las
variables EOD de ingreso. En consecuencia, se considera una sensibilidad útil para diagnóstico, pero
no una variable central del modelo de reporte.

Como sensibilidad computacional y de robustez, también se estimaron modelos equivalentes en Larch con
muestras mayores. En particular, se probaron muestras de 5\% y 10\%, aprovechando que Larch permitió
trabajar con mayor volumen de datos sin los problemas de memoria observados previamente. Los
resultados mantuvieron métricas de ajuste y patrones generales consistentes con la estimación en
Biogeme al 2\%. Esto entrega una confirmación adicional de que las señales principales no dependen
únicamente del submuestreo inicial.

En síntesis, el modelo territorial de trabajo queda definido por una estructura común con variables
de viaje, temporalidad, Censo, OSM y EOD 2012. La decisión metodológica pendiente es seleccionar qué
representación del bloque EOD resulta más adecuada para la tesis: una lectura continua de gradientes
territoriales de ingreso o una lectura discreta basada en zonas de alta concentración relativa. En
ambos casos, la especificación debe entenderse como un modelo principal de trabajo, no como el cierre
definitivo de la investigación.

Aún queda pendiente evaluar otras familias de variables que podrían complementar la lectura
territorial, especialmente aquellas vinculadas con calidad de infraestructura de paraderos y
estaciones, accesibilidad peatonal, seguridad del entorno, disponibilidad de puntos de carga o pago,
y medidas operacionales de confiabilidad, regularidad y ocupación. Estas variables podrían mejorar
la capacidad explicativa del modelo y, sobre todo, aumentar su valor interpretativo para una
autoridad de transporte.

\appendix

\section{Tablas completas de parámetros}
\label{sec:anexo_parametros_completos}

Esta sección presenta las tablas extendidas de parámetros, incluyendo errores estándar robustos, para
las dos especificaciones candidatas del bloque EOD 2012.

\begin{landscape}
\scriptsize
\setlength{\tabcolsep}{2pt}
\begin{longtable}{p{0.18\linewidth}rrrrrrrrr}
\caption{Tabla extendida de parámetros del modelo con proxy EOD continuo}
\label{tab:parametros_extendida_eod_continuo}\\
\toprule
Variable &
\texttt{BIP} Est. & \texttt{BIP} s.e. & \texttt{BIP} t &
\texttt{QRO} Est. & \texttt{QRO} s.e. & \texttt{QRO} t &
\texttt{QRR} Est. & \texttt{QRR} s.e. & \texttt{QRR} t \\
\midrule
\endfirsthead

\toprule
Variable &
\texttt{BIP} Est. & \texttt{BIP} s.e. & \texttt{BIP} t &
\texttt{QRO} Est. & \texttt{QRO} s.e. & \texttt{QRO} t &
\texttt{QRR} Est. & \texttt{QRR} s.e. & \texttt{QRR} t \\
\midrule
\endhead

\bottomrule
\endfoot

\texttt{ASC} & -- & -- & -- & -1.9172 & 0.0822 & -23.32 & -3.5839 & 0.1722 & -20.81 \\
\texttt{ANIO\_2025} & 0 & -- & -- & 0.2965 & 0.008898 & 33.32 & 0.2332 & 0.0183 & 12.72 \\
\texttt{LAB\_PM} & 0 & -- & -- & 0.0976 & 0.0144 & 6.76 & 0.4604 & 0.0302 & 15.25 \\
\texttt{LAB\_PT} & 0 & -- & -- & 0.1571 & 0.0177 & 8.87 & 0.2396 & 0.0348 & 6.89 \\
\texttt{NO\_LAB} & 0 & -- & -- & 0.1624 & 0.0176 & 9.23 & 0.0510 & 0.0378 & 1.35 \\
\texttt{LOG\_DEMAND} & 0 & -- & -- & -0.0344 & 0.007988 & -4.31 & -0.0671 & 0.0158 & -4.24 \\
\texttt{Bus stops} & 0 & -- & -- & 0.000826 & 0.0138 & 0.06 & -0.0910 & 0.0312 & -2.92 \\
\texttt{Bus lines} & 0 & -- & -- & 0.0164 & 0.0146 & 1.13 & 0.0901 & 0.0310 & 2.91 \\
\texttt{Metro lines} & 0 & -- & -- & 0.0488 & 0.0204 & 2.39 & 0.1235 & 0.0393 & 3.14 \\
\texttt{T\_VEH} & 0.000020 & 0.000030 & 0.66 & 0.000071 & 0.000030 & 2.39 & -0.000006 & 0.000030 & -0.19 \\
\texttt{T\_ESPERA\_INI} & -0.001706 & 0.000058 & -29.21 & -0.001729 & 0.000059 & -29.39 & -0.001408 & 0.000070 & -20.08 \\
\texttt{T\_ESPERA\_TRASB} & -0.000278 & 0.000077 & -3.59 & -0.000445 & 0.000079 & -5.60 & 0.000045 & 0.000075 & 0.60 \\
\texttt{N\_TRASB} & -0.0597 & 0.0397 & -1.50 & -0.1248 & 0.0396 & -3.15 & -0.1556 & 0.0407 & -3.83 \\
\texttt{Educación univ. 18+} & 0 & -- & -- & 0.0815 & 0.009128 & 8.93 & 0.3073 & 0.0187 & 16.47 \\
\texttt{Hacinamiento} & 0 & -- & -- & -0.0652 & 0.0115 & -5.66 & -0.0492 & 0.0275 & -1.79 \\
\texttt{Discapacidad} & 0 & -- & -- & 0.002760 & 0.0115 & 0.24 & -0.0988 & 0.0269 & -3.67 \\
\texttt{Inmigrantes} & 0 & -- & -- & 0.1007 & 0.0102 & 9.91 & 0.0865 & 0.0235 & 3.68 \\
\texttt{Mujeres} & 0 & -- & -- & -0.000229 & 0.003522 & -0.06 & 0.0212 & 0.008919 & 2.37 \\
\texttt{Asistencia parvularia} & 0 & -- & -- & 0.005331 & 0.007041 & 0.76 & 0.0487 & 0.0156 & 3.11 \\
\texttt{Centros deportivos OSM} & 0 & -- & -- & -0.009242 & 0.004398 & -2.10 & 0.007836 & 0.008764 & 0.89 \\
\texttt{Convenience OSM} & 0 & -- & -- & 0.0135 & 0.004399 & 3.07 & 0.0414 & 0.008815 & 4.69 \\
\texttt{Playgrounds OSM} & 0 & -- & -- & 0.0283 & 0.005005 & 5.65 & 0.0662 & 0.0103 & 6.44 \\
\texttt{Escuelas OSM} & 0 & -- & -- & -0.000137 & 0.005063 & -0.03 & -0.0284 & 0.0112 & -2.54 \\
\texttt{Universidades OSM} & 0 & -- & -- & -0.0197 & 0.002153 & -9.15 & -0.0244 & 0.004556 & -5.35 \\
\texttt{Refugios transporte OSM} & 0 & -- & -- & -0.0306 & 0.006576 & -4.65 & 0.0375 & 0.0143 & 2.62 \\
\texttt{Accesos Metro OSM} & 0 & -- & -- & -0.009577 & 0.003380 & -2.83 & -0.0370 & 0.007011 & -5.27 \\
\texttt{ABC1\_proxy EOD} & 0 & -- & -- & -0.0154 & 0.008262 & -1.87 & -0.005278 & 0.0157 & -0.34 \\
\texttt{D+E\_proxy EOD} & 0 & -- & -- & 0.0369 & 0.008138 & 4.53 & -0.0123 & 0.0183 & -0.67 \\

\end{longtable}
\end{landscape}

\begin{landscape}
\scriptsize
\setlength{\tabcolsep}{2pt}
\begin{longtable}{p{0.18\linewidth}rrrrrrrrr}
\caption{Tabla extendida de parámetros del modelo con proxy EOD dummy}
\label{tab:parametros_extendida_eod_dummy}\\
\toprule
Variable &
\texttt{BIP} Est. & \texttt{BIP} s.e. & \texttt{BIP} t &
\texttt{QRO} Est. & \texttt{QRO} s.e. & \texttt{QRO} t &
\texttt{QRR} Est. & \texttt{QRR} s.e. & \texttt{QRR} t \\
\midrule
\endfirsthead

\toprule
Variable &
\texttt{BIP} Est. & \texttt{BIP} s.e. & \texttt{BIP} t &
\texttt{QRO} Est. & \texttt{QRO} s.e. & \texttt{QRO} t &
\texttt{QRR} Est. & \texttt{QRR} s.e. & \texttt{QRR} t \\
\midrule
\endhead

\bottomrule
\endfoot

\texttt{ASC} & -- & -- & -- & -1.9081 & 0.0824 & -23.16 & -3.5086 & 0.1714 & -20.47 \\
\texttt{ANIO\_2025} & 0 & -- & -- & 0.2963 & 0.008897 & 33.30 & 0.2329 & 0.0183 & 12.70 \\
\texttt{LAB\_PM} & 0 & -- & -- & 0.0955 & 0.0145 & 6.60 & 0.4529 & 0.0303 & 14.97 \\
\texttt{LAB\_PT} & 0 & -- & -- & 0.1522 & 0.0177 & 8.60 & 0.2305 & 0.0348 & 6.62 \\
\texttt{NO\_LAB} & 0 & -- & -- & 0.1603 & 0.0176 & 9.11 & 0.0447 & 0.0378 & 1.18 \\
\texttt{LOG\_DEMAND} & 0 & -- & -- & -0.0370 & 0.007954 & -4.65 & -0.0722 & 0.0158 & -4.56 \\
\texttt{Bus stops} & 0 & -- & -- & -0.002211 & 0.0140 & -0.16 & -0.1134 & 0.0315 & -3.60 \\
\texttt{Bus lines} & 0 & -- & -- & 0.0251 & 0.0146 & 1.72 & 0.1023 & 0.0310 & 3.29 \\
\texttt{Metro lines} & 0 & -- & -- & 0.0571 & 0.0204 & 2.80 & 0.1251 & 0.0392 & 3.19 \\
\texttt{T\_VEH} & 0.000020 & 0.000030 & 0.67 & 0.000069 & 0.000030 & 2.32 & -0.000008 & 0.000030 & -0.25 \\
\texttt{T\_ESPERA\_INI} & -0.001714 & 0.000058 & -29.34 & -0.001732 & 0.000059 & -29.43 & -0.001406 & 0.000070 & -20.04 \\
\texttt{T\_ESPERA\_TRASB} & -0.000277 & 0.000077 & -3.58 & -0.000443 & 0.000079 & -5.58 & 0.000046 & 0.000075 & 0.61 \\
\texttt{N\_TRASB} & -0.0644 & 0.0398 & -1.62 & -0.1285 & 0.0396 & -3.24 & -0.1572 & 0.0407 & -3.86 \\
\texttt{Educación univ. 18+} & 0 & -- & -- & 0.0784 & 0.009077 & 8.63 & 0.3194 & 0.0186 & 17.16 \\
\texttt{Hacinamiento} & 0 & -- & -- & -0.0556 & 0.0115 & -4.84 & -0.0520 & 0.0276 & -1.88 \\
\texttt{Discapacidad} & 0 & -- & -- & 0.0123 & 0.0113 & 1.09 & -0.1057 & 0.0263 & -4.02 \\
\texttt{Inmigrantes} & 0 & -- & -- & 0.1033 & 0.0101 & 10.24 & 0.0962 & 0.0235 & 4.09 \\
\texttt{Mujeres} & 0 & -- & -- & 0.000530 & 0.003571 & 0.15 & 0.0241 & 0.008994 & 2.68 \\
\texttt{Asistencia parvularia} & 0 & -- & -- & 0.009437 & 0.007118 & 1.33 & 0.0672 & 0.0161 & 4.17 \\
\texttt{Centros deportivos OSM} & 0 & -- & -- & -0.008577 & 0.004410 & -1.94 & 0.0113 & 0.008755 & 1.29 \\
\texttt{Convenience OSM} & 0 & -- & -- & 0.0119 & 0.004338 & 2.74 & 0.0429 & 0.008681 & 4.94 \\
\texttt{Playgrounds OSM} & 0 & -- & -- & 0.0283 & 0.005022 & 5.63 & 0.0644 & 0.0103 & 6.26 \\
\texttt{Escuelas OSM} & 0 & -- & -- & 0.001978 & 0.005052 & 0.39 & -0.0262 & 0.0112 & -2.35 \\
\texttt{Universidades OSM} & 0 & -- & -- & -0.0178 & 0.002213 & -8.03 & -0.0211 & 0.004714 & -4.47 \\
\texttt{Refugios transporte OSM} & 0 & -- & -- & -0.0304 & 0.006650 & -4.58 & 0.0328 & 0.0144 & 2.28 \\
\texttt{Accesos Metro OSM} & 0 & -- & -- & -0.0104 & 0.003385 & -3.06 & -0.0340 & 0.006940 & -4.90 \\
\texttt{Alta conc. ABC1\_proxy} & 0 & -- & -- & -0.0264 & 0.006044 & -4.38 & -0.0421 & 0.0122 & -3.44 \\
\texttt{Alta conc. D+E\_proxy} & 0 & -- & -- & -0.002550 & 0.006913 & -0.37 & -0.0110 & 0.0171 & -0.65 \\

\end{longtable}
\end{landscape}

\section{Odds ratios del modelo principal}
\label{sec:anexo_odds_ratios}

Esta sección presenta los \emph{odds ratios} de las dos especificaciones candidatas. Para variables
comunes al viaje o a la zona, \texttt{BIP} se omite porque corresponde a la alternativa de
referencia. Para variables alternativa-específicas, como tiempos y transbordos, se reportan también
los valores de \texttt{BIP}, ya que esos coeficientes son estimados directamente. La columna de
porcentaje se calcula como \((OR-1)\times 100\), por lo que representa cambio porcentual en
\emph{odds}, no en probabilidad.

\begin{landscape}
\scriptsize
\setlength{\tabcolsep}{2pt}
\begin{longtable}{p{0.18\linewidth}p{0.12\linewidth}rrrrrrrrr}
\caption{Odds ratios del modelo con proxy EOD continuo}
\label{tab:odds_ratios_eod_continuo}\\
\toprule
Variable &
Unidad OR &
\texttt{BIP} OR & \texttt{BIP} \% & \texttt{BIP} t &
\texttt{QRO} OR & \texttt{QRO} \% & \texttt{QRO} t &
\texttt{QRR} OR & \texttt{QRR} \% & \texttt{QRR} t \\
\midrule
\endfirsthead

\toprule
Variable &
Unidad OR &
\texttt{BIP} OR & \texttt{BIP} \% & \texttt{BIP} t &
\texttt{QRO} OR & \texttt{QRO} \% & \texttt{QRO} t &
\texttt{QRR} OR & \texttt{QRR} \% & \texttt{QRR} t \\
\midrule
\endhead

\bottomrule
\endfoot

\texttt{ASC} & Constante & -- & -- & -- & 0.1470 & -85.30 & -23.32 & 0.0278 & -97.22 & -20.81 \\
\texttt{ANIO\_2025} & Dummy 0--1 & -- & -- & -- & 1.3451 & 34.51 & 33.32 & 1.2626 & 26.26 & 12.72 \\
\texttt{LAB\_PM} & Dummy 0--1 & -- & -- & -- & 1.1025 & 10.25 & 6.76 & 1.5847 & 58.47 & 15.25 \\
\texttt{LAB\_PT} & Dummy 0--1 & -- & -- & -- & 1.1701 & 17.01 & 8.87 & 1.2707 & 27.07 & 6.89 \\
\texttt{NO\_LAB} & Dummy 0--1 & -- & -- & -- & 1.1763 & 17.63 & 9.23 & 1.0523 & 5.23 & 1.35 \\
\texttt{LOG\_DEMAND} & 1 unidad log & -- & -- & -- & 0.9662 & -3.38 & -4.31 & 0.9351 & -6.49 & -4.24 \\
\texttt{Bus stops} & 1 unidad log & -- & -- & -- & 1.0008 & 0.08 & 0.06 & 0.9130 & -8.70 & -2.92 \\
\texttt{Bus lines} & 1 unidad log & -- & -- & -- & 1.0166 & 1.66 & 1.13 & 1.0943 & 9.43 & 2.91 \\
\texttt{Metro lines} & 1 unidad log & -- & -- & -- & 1.0501 & 5.01 & 2.39 & 1.1315 & 13.15 & 3.14 \\
\texttt{T\_VEH} & 1 minuto & 1.0012 & 0.12 & 0.66 & 1.0043 & 0.43 & 2.39 & 0.9997 & -0.03 & -0.19 \\
\texttt{T\_ESPERA\_INI} & 1 minuto & 0.9027 & -9.73 & -29.21 & 0.9014 & -9.86 & -29.39 & 0.9190 & -8.10 & -20.08 \\
\texttt{T\_ESPERA\_TRASB} & 1 minuto & 0.9835 & -1.65 & -3.59 & 0.9737 & -2.63 & -5.60 & 1.0027 & 0.27 & 0.60 \\
\texttt{N\_TRASB} & 1 transbordo & 0.9420 & -5.80 & -1.50 & 0.8827 & -11.73 & -3.15 & 0.8559 & -14.41 & -3.83 \\
\texttt{Educación univ. 18+} & +1 d.e. & -- & -- & -- & 1.0849 & 8.49 & 8.93 & 1.3597 & 35.97 & 16.47 \\
\texttt{Hacinamiento} & +1 d.e. & -- & -- & -- & 0.9369 & -6.31 & -5.66 & 0.9519 & -4.81 & -1.79 \\
\texttt{Discapacidad} & +1 d.e. & -- & -- & -- & 1.0028 & 0.28 & 0.24 & 0.9059 & -9.41 & -3.67 \\
\texttt{Inmigrantes} & +1 d.e. & -- & -- & -- & 1.1059 & 10.59 & 9.91 & 1.0904 & 9.04 & 3.68 \\
\texttt{Mujeres} & +1 d.e. & -- & -- & -- & 0.9998 & -0.02 & -0.06 & 1.0214 & 2.14 & 2.37 \\
\texttt{Asistencia parvularia} & +1 d.e. & -- & -- & -- & 1.0053 & 0.53 & 0.76 & 1.0499 & 4.99 & 3.11 \\
\texttt{Centros deportivos OSM} & +1 d.e. & -- & -- & -- & 0.9908 & -0.92 & -2.10 & 1.0079 & 0.79 & 0.89 \\
\texttt{Convenience OSM} & +1 d.e. & -- & -- & -- & 1.0136 & 1.36 & 3.07 & 1.0422 & 4.22 & 4.69 \\
\texttt{Playgrounds OSM} & +1 d.e. & -- & -- & -- & 1.0287 & 2.87 & 5.65 & 1.0684 & 6.84 & 6.44 \\
\texttt{Escuelas OSM} & +1 d.e. & -- & -- & -- & 0.9999 & -0.01 & -0.03 & 0.9720 & -2.80 & -2.54 \\
\texttt{Universidades OSM} & +1 d.e. & -- & -- & -- & 0.9805 & -1.95 & -9.15 & 0.9759 & -2.41 & -5.35 \\
\texttt{Refugios transporte OSM} & +1 d.e. & -- & -- & -- & 0.9699 & -3.01 & -4.65 & 1.0382 & 3.82 & 2.62 \\
\texttt{Accesos Metro OSM} & +1 d.e. & -- & -- & -- & 0.9905 & -0.95 & -2.83 & 0.9637 & -3.63 & -5.27 \\
\texttt{ABC1\_proxy EOD} & +1 d.e. & -- & -- & -- & 0.9847 & -1.53 & -1.87 & 0.9947 & -0.53 & -0.34 \\
\texttt{D+E\_proxy EOD} & +1 d.e. & -- & -- & -- & 1.0376 & 3.76 & 4.53 & 0.9878 & -1.22 & -0.67 \\

\end{longtable}
\end{landscape}

\begin{landscape}
\scriptsize
\setlength{\tabcolsep}{2pt}
\begin{longtable}{p{0.18\linewidth}p{0.12\linewidth}rrrrrrrrr}
\caption{Odds ratios del modelo con proxy EOD dummy}
\label{tab:odds_ratios_eod_dummy}\\
\toprule
Variable &
Unidad OR &
\texttt{BIP} OR & \texttt{BIP} \% & \texttt{BIP} t &
\texttt{QRO} OR & \texttt{QRO} \% & \texttt{QRO} t &
\texttt{QRR} OR & \texttt{QRR} \% & \texttt{QRR} t \\
\midrule
\endfirsthead

\toprule
Variable &
Unidad OR &
\texttt{BIP} OR & \texttt{BIP} \% & \texttt{BIP} t &
\texttt{QRO} OR & \texttt{QRO} \% & \texttt{QRO} t &
\texttt{QRR} OR & \texttt{QRR} \% & \texttt{QRR} t \\
\midrule
\endhead

\bottomrule
\endfoot

\texttt{ASC} & Constante & -- & -- & -- & 0.1484 & -85.16 & -23.16 & 0.0299 & -97.01 & -20.47 \\
\texttt{ANIO\_2025} & Dummy 0--1 & -- & -- & -- & 1.3448 & 34.48 & 33.30 & 1.2623 & 26.23 & 12.70 \\
\texttt{LAB\_PM} & Dummy 0--1 & -- & -- & -- & 1.1002 & 10.02 & 6.60 & 1.5729 & 57.29 & 14.97 \\
\texttt{LAB\_PT} & Dummy 0--1 & -- & -- & -- & 1.1643 & 16.43 & 8.60 & 1.2592 & 25.92 & 6.62 \\
\texttt{NO\_LAB} & Dummy 0--1 & -- & -- & -- & 1.1739 & 17.39 & 9.11 & 1.0457 & 4.57 & 1.18 \\
\texttt{LOG\_DEMAND} & 1 unidad log & -- & -- & -- & 0.9637 & -3.63 & -4.65 & 0.9303 & -6.97 & -4.56 \\
\texttt{Bus stops} & 1 unidad log & -- & -- & -- & 0.9978 & -0.22 & -0.16 & 0.8928 & -10.72 & -3.60 \\
\texttt{Bus lines} & 1 unidad log & -- & -- & -- & 1.0254 & 2.54 & 1.72 & 1.1077 & 10.77 & 3.29 \\
\texttt{Metro lines} & 1 unidad log & -- & -- & -- & 1.0587 & 5.87 & 2.80 & 1.1333 & 13.33 & 3.19 \\
\texttt{T\_VEH} & 1 minuto & 1.0012 & 0.12 & 0.67 & 1.0042 & 0.42 & 2.32 & 0.9995 & -0.05 & -0.25 \\
\texttt{T\_ESPERA\_INI} & 1 minuto & 0.9023 & -9.77 & -29.34 & 0.9013 & -9.87 & -29.43 & 0.9191 & -8.09 & -20.04 \\
\texttt{T\_ESPERA\_TRASB} & 1 minuto & 0.9835 & -1.65 & -3.58 & 0.9738 & -2.62 & -5.58 & 1.0027 & 0.27 & 0.61 \\
\texttt{N\_TRASB} & 1 transbordo & 0.9376 & -6.24 & -1.62 & 0.8794 & -12.06 & -3.24 & 0.8546 & -14.54 & -3.86 \\
\texttt{Educación univ. 18+} & +1 d.e. & -- & -- & -- & 1.0815 & 8.15 & 8.63 & 1.3763 & 37.63 & 17.16 \\
\texttt{Hacinamiento} & +1 d.e. & -- & -- & -- & 0.9459 & -5.41 & -4.84 & 0.9494 & -5.06 & -1.88 \\
\texttt{Discapacidad} & +1 d.e. & -- & -- & -- & 1.0124 & 1.24 & 1.09 & 0.8997 & -10.03 & -4.02 \\
\texttt{Inmigrantes} & +1 d.e. & -- & -- & -- & 1.1088 & 10.88 & 10.24 & 1.1010 & 10.10 & 4.09 \\
\texttt{Mujeres} & +1 d.e. & -- & -- & -- & 1.0005 & 0.05 & 0.15 & 1.0244 & 2.44 & 2.68 \\
\texttt{Asistencia parvularia} & +1 d.e. & -- & -- & -- & 1.0095 & 0.95 & 1.33 & 1.0695 & 6.95 & 4.17 \\
\texttt{Centros deportivos OSM} & +1 d.e. & -- & -- & -- & 0.9915 & -0.85 & -1.94 & 1.0113 & 1.13 & 1.29 \\
\texttt{Convenience OSM} & +1 d.e. & -- & -- & -- & 1.0119 & 1.19 & 2.74 & 1.0438 & 4.38 & 4.94 \\
\texttt{Playgrounds OSM} & +1 d.e. & -- & -- & -- & 1.0287 & 2.87 & 5.63 & 1.0665 & 6.65 & 6.26 \\
\texttt{Escuelas OSM} & +1 d.e. & -- & -- & -- & 1.0020 & 0.20 & 0.39 & 0.9741 & -2.59 & -2.35 \\
\texttt{Universidades OSM} & +1 d.e. & -- & -- & -- & 0.9824 & -1.76 & -8.03 & 0.9792 & -2.08 & -4.47 \\
\texttt{Refugios transporte OSM} & +1 d.e. & -- & -- & -- & 0.9700 & -3.00 & -4.58 & 1.0334 & 3.34 & 2.28 \\
\texttt{Accesos Metro OSM} & +1 d.e. & -- & -- & -- & 0.9897 & -1.03 & -3.06 & 0.9665 & -3.35 & -4.90 \\
\texttt{Alta conc. ABC1\_proxy} & 1 d.e. dummy & -- & -- & -- & 0.9739 & -2.61 & -4.38 & 0.9588 & -4.12 & -3.44 \\
\texttt{Alta conc. D+E\_proxy} & 1 d.e. dummy & -- & -- & -- & 0.9975 & -0.25 & -0.37 & 0.9890 & -1.10 & -0.65 \\

\end{longtable}
\end{landscape}


\end{document}
